\section{General algebraic extensions}

Throughout this section $S=(K,(\Omega,\mathcal A,\nu),\phi)$ is an
adelic structure and $L/K$ is algebraic. We follow
\cite[\S3.4.2]{chen-moriwaki}. A finite-stage parameter space $\Omega_F$
is identified with the absolute values of $F$ extending the specified base
absolute value, using the unique extension across the purely inseparable
part. Thus restriction has its literal meaning also at inseparable stages.

\begin{definition}[Finite stages]\label{def:finite-stages}
\source{chen-moriwaki-adelic:page-0207}
\lean{ArakelovGeometry.AdelicStructure.FiniteIntermediateField,ArakelovGeometry.AdelicStructure.FiniteIntermediateField.bot,ArakelovGeometry.AdelicStructure.FiniteIntermediateField.sup,ArakelovGeometry.AdelicStructure.instCoeFiniteIntermediateFieldIntermediateField,ArakelovGeometry.AdelicStructure.instCoeSortFiniteIntermediateFieldType}
Let $\mathscr E_{L/K}$ be the intermediate fields $F$ with $[F:K]<\infty$,
ordered by inclusion. Its least element is $K$, and the compositum $FG$
is an upper bound for $F,G\in\mathscr E_{L/K}$.
\end{definition}
\begin{proof}
The field $K$ has degree one over itself. If $f_1,\ldots,f_r$ span $F/K$
and $g_1,\ldots,g_s$ span $G/K$, their products span a finite-dimensional
$K$-subalgebra containing $F$ and $G$. This subalgebra is a field:
multiplication by a nonzero element is an injective linear endomorphism,
and hence is surjective and has $1$ in its image. It is therefore $FG$.
\end{proof}

\begin{lemma}[Directedness and exhaustion]\label{lem:finite-stage-exhaustion}
\uses{def:finite-stages}
\source{chen-moriwaki-adelic:page-0207,chen-moriwaki-adelic:page-0211}
\lean{ArakelovGeometry.AdelicStructure.FiniteIntermediateField.left_le_sup,ArakelovGeometry.AdelicStructure.FiniteIntermediateField.right_le_sup,ArakelovGeometry.AdelicStructure.Algebraic.ofElement,ArakelovGeometry.AdelicStructure.Algebraic.mem_ofElement}
The inclusions $F\subseteq FG$ and $G\subseteq FG$ hold. Every $a\in L$
belongs to the finite stage $K(a)$; consequently every finite subset of $L$
is contained in one finite stage.
\end{lemma}
\begin{proof}
The compositum contains its two generating subfields. If $a$ has minimal
polynomial of degree $d$, then $1,a,\ldots,a^{d-1}$ span $K(a)$, and $a$
belongs to this field. Taking successive composita handles any finite set.
\end{proof}

\begin{lemma}[Extending algebraic embeddings]\label{lem:algebraic-embedding-extension}
\dcref{IsAlgClosed.lift}
Let $E/F$ be algebraic and let $j:F\hookrightarrow C$ be a field embedding
with $C$ algebraically closed. Then $j$ extends to an embedding $E\hookrightarrow C$.
\end{lemma}
\begin{proof}
Order embeddings of intermediate fields into $C$ extending $j$ by
extension. The union of a chain is an embedding of the union of its domains,
so Zorn's lemma gives a maximal embedding $j_D:D\hookrightarrow C$.
If $a\in E\setminus D$, its minimal polynomial over $D$, transported by
$j_D$, has a root $b\in C$. Evaluation at $b$ induces a homomorphism
$D[a]\to C$ extending $j_D$. The kernel in $D[T]$ is generated by the
irreducible minimal polynomial, so this is an embedding of $D(a)$.
This contradicts maximality. Hence $D=E$.
\end{proof}

\begin{lemma}[Extension of an absolute value to an algebraic field extension]
\label{lem:algebraic-absolute-value-extension}
\leanok
\uses{def:absolute-value-restriction,lem:complete-extension,lem:completion-embedding,lem:algebraic-embedding-extension}
\source{chen-moriwaki-adelic:page-0208}
\lean{AbsoluteValue.exists_extension}
Every real absolute value $v$ on a field $F$ extends to every algebraic
field extension $E/F$.
\end{lemma}
\begin{proof}
Embed $F$ isometrically into its completion $F_v$, and let $C$ be an
algebraic closure of $F_v$. Extend the composite embedding $F\to C$
to an embedding $E\to C$ by Lemma~\ref{lem:algebraic-embedding-extension}.
The completed absolute value has a unique extension to $C$ by
Lemma~\ref{lem:complete-extension}. Pulling it back along $E\to C$
gives an absolute value whose restriction to $F$ is $v$.
\end{proof}

\begin{proposition}[Surjectivity of restriction]\label{prop:algebraic-restriction-surjective}
\leanok
\uses{def:extension-parameters,def:restriction-parameters,lem:algebraic-absolute-value-extension,lem:extension-coordinates}
\source{chen-moriwaki-adelic:page-0208}
\lean{ArakelovGeometry.AdelicStructure.Algebraic.surjective_restrict,ArakelovGeometry.AdelicStructure.Algebraic.surjective_proj}
For every intermediate field $F$ of $L/K$, restriction $\Omega_L\to\Omega_F$
is surjective. In particular $\pi_{L/K}:\Omega_L\to\Omega$ is surjective.
\end{proposition}
\begin{proof}
Fix $(\omega,v)\in\Omega_F$. By
Lemma~\ref{lem:algebraic-absolute-value-extension}, $v$ extends to an
absolute value $w$ on $L$. Its restriction to $K$ equals the restriction
of $v$, namely $\phi_\omega$. Thus $(\omega,w)$ is a point of
$\Omega_L$ restricting to $(\omega,v)$. For the base projection extend
$\phi_\omega$ directly from $K$ to $L$.
\end{proof}

\begin{definition}[Finite projections and fibres]\label{def:algebraic-projections}
\uses{def:extension-parameters,def:restriction-parameters,def:finite-stages,def:finite-extension}
\source{chen-moriwaki-adelic:page-0208}
\lean{ArakelovGeometry.AdelicStructure.Algebraic.projFinite,ArakelovGeometry.AdelicStructure.Algebraic.FiniteBaseFiber,ArakelovGeometry.AdelicStructure.Algebraic.BaseFiber,ArakelovGeometry.AdelicStructure.Algebraic.projFiniteFiber,ArakelovGeometry.AdelicStructure.Algebraic.finiteFiber}
For $F\in\mathscr E_{L/K}$ write $p_F:\Omega_L\to\Omega_F$ for
restriction. Set $X_\omega=\pi_{L/K}^{-1}\{\omega\}$ and
$X_{F,\omega}=\pi_{F/K}^{-1}\{\omega\}$, and denote the restricted map
by $p_{F,\omega}:X_\omega\to X_{F,\omega}$. For $x\in\Omega_F$ put
$M_{L,x}=p_F^{-1}\{x\}$. In the separable-part realization of $\Omega_F$,
$p_F$ restricts first to $F$ and then to the separable closure of $K$ in $F$.
\end{definition}

\begin{lemma}[Transitions between finite stages]\label{lem:finite-stage-transitions}
\leanok
\uses{def:algebraic-projections,lem:finite-separable-part,def:absolute-value-restriction}
\source{chen-moriwaki-adelic:page-0208}
\lean{ArakelovGeometry.AdelicStructure.Algebraic.finiteSeparableInclusion,ArakelovGeometry.AdelicStructure.Algebraic.projFiniteOfLE,ArakelovGeometry.AdelicStructure.Algebraic.projFiniteOfLE_projFinite,ArakelovGeometry.AdelicStructure.Algebraic.projFiniteFiberOfLE,ArakelovGeometry.AdelicStructure.Algebraic.projFiniteFiberOfLE_projFiniteFiber}
For finite intermediate fields $F\subseteq G\subseteq L$, inclusion
induces $F^{\mathrm{sep}}\hookrightarrow G^{\mathrm{sep}}$.
Restriction of absolute values along this inclusion defines
$p_{G/F}:\Omega_G\to\Omega_F$, preserves the base parameter,
and satisfies $p_{G/F}\circ p_G=p_F$. The induced maps on each
fibre satisfy the same compatibility.
\end{lemma}
\begin{proof}
An element separable over $K$ remains separable under a $K$-embedding:
the embedding transports its minimal polynomial and preserves its
distinct roots. Thus inclusion maps the separable part of $F$ into
that of $G$. Pulling back an absolute value gives an absolute value
whose restriction to $K$ is unchanged. For $a\in F^{\mathrm{sep}}$,
restriction first to $G^{\mathrm{sep}}$ and then to $F^{\mathrm{sep}}$
evaluates the original absolute value at the same element $a$.
The base parameter is unchanged, so the parameter maps agree and
restrict to the asserted maps on each fibre.
\end{proof}

\begin{lemma}[Finite fibres are discrete]\label{lem:finite-stage-discrete}
\leanok
\uses{def:algebraic-projections,lem:finite-fibre-sum,lem:fibre-separation}
\source{chen-moriwaki-adelic:page-0208}
\lean{ArakelovGeometry.AdelicStructure.Algebraic.measurableSet_finiteBaseFiber}
The set $X_{F,\omega}$ is finite and every subset is measurable for its
induced measurable structure.
\end{lemma}
\begin{proof}
Finiteness follows from the finite separable description of the fibre.
For a point $x$ and each other point $y$, measurable separation gives a
measurable set containing $x$ but not $y$; the intersection of these
finitely many sets is $\{x\}$. Every subset is a finite union of singletons.
\end{proof}

\begin{definition}[Projective measurable structures]\label{def:projective-measurable}
\uses{def:algebraic-projections,lem:finite-stage-discrete}
\source{chen-moriwaki-adelic:page-0208}
\lean{ArakelovGeometry.AdelicStructure.Algebraic.projectiveMeasurableSpace,ArakelovGeometry.AdelicStructure.Algebraic.fibreMeasurableSpace}
Let $\mathcal A_L$ be the smallest $\sigma$-algebra on $\Omega_L$ making
every $p_F$ measurable. Let $\mathcal A_{L,\omega}$ be the smallest
$\sigma$-algebra on $X_\omega$ making every $p_{F,\omega}$ measurable,
where the finite fibres are discrete measurable spaces.
\end{definition}

\begin{lemma}[Projective measurability of projections and evaluations]\label{lem:projective-evaluations}
\leanok
\uses{def:projective-measurable,lem:finite-stage-exhaustion,lem:finite-measurable-space}
\source{chen-moriwaki-adelic:page-0211,chen-moriwaki-adelic:page-0215}
\lean{ArakelovGeometry.AdelicStructure.Algebraic.measurable_projFinite_projective,ArakelovGeometry.AdelicStructure.Algebraic.measurable_proj_projective,ArakelovGeometry.AdelicStructure.Algebraic.measurable_eval_projective,ArakelovGeometry.AdelicStructure.Algebraic.measurable_log_projective}
Each $p_F$ is measurable for $\mathcal A_L$. The base projection
$\pi_{L/K}$ and every evaluation $x\mapsto |a|_x$, $a\in L$, are
$\mathcal A_L$-measurable. So is $x\mapsto\log|a|_x$, with the convention
$\log0=0$.
\end{lemma}
\begin{proof}
The first assertion is the defining property. The base projection factors
through the stage $K$. For $F=K(a)$, evaluation at $a$ is the composite
of $p_F$ and the measurable finite-stage evaluation. Composing with the
Borel function $\log$ proves the last assertion; its assigned value at
zero does not affect Borel measurability.
\end{proof}

\begin{proposition}[Intrinsic description of the measurable structure]\label{prop:projective-measurable-equality}
\leanok
\uses{lem:projective-evaluations,def:extension-parameters,lem:measurable-restriction,lem:finite-measurable-space}
\source{chen-moriwaki-adelic:page-0214,chen-moriwaki-adelic:page-0215}
\lean{ArakelovGeometry.AdelicStructure.Algebraic.projectiveMeasurableSpace_eq}
$\mathcal A_L$ equals the $\sigma$-algebra generated by $\pi_{L/K}$ and
all evaluations at elements of $L$.
\end{proposition}
\begin{proof}
Lemma~\ref{lem:projective-evaluations} gives one inclusion. For the other,
fix $F$. The finite-stage $\sigma$-algebra is generated by its base
projection and its evaluations. Their composites with $p_F$ are the
base projection of $\Omega_L$ and the evaluations at elements of $F$.
Thus $p_F$ is measurable for the claimed $\sigma$-algebra. This holds
for every $F$, which proves the reverse inclusion.
\end{proof}

\begin{lemma}[Compactness of extending absolute values by evaluation]
\label{lem:compact-absolute-value-extensions}
\leanok
\lean{AbsoluteValue.isClosed_range_extensions,AbsoluteValue.finite_range_extensions_apply,AbsoluteValue.isCompact_range_extensions}
\source{chen-moriwaki-adelic:page-0194,chen-moriwaki-adelic:page-0208,chen-moriwaki-adelic:page-0209}
\uses{lem:finite-values-via-embeddings,lem:finite-stage-exhaustion}
Let $L/K$ be algebraic and let $v$ be an absolute value on $K$.
The set of functions $L\to\mathbb R$ arising from absolute values
on $L$ extending $v$ is compact in the topology of pointwise convergence.
For each $a\in L$, the set of possible values $w(a)$ is finite.
\end{lemma}
\begin{proof}
Put $E=K(a)$. Restriction to $E$ shows that each value $w(a)$ comes
from an absolute value on $E$ extending $v$. There are finitely many
such values, because $E/K$ is finite and the extending absolute values
form a finite set. Denote this finite subset of $\mathbb R$ by $T_a$.

The set in question is closed in $\mathbb R^L$. Indeed, it is exactly
the set of functions $f$ satisfying
\[
 \begin{gathered}
 f(0)=0,\quad f(1)=1,\quad f(a)\geq0,\quad f(ab)=f(a)f(b),\\
 f(a+b)\leq f(a)+f(b),\qquad f(c)=v(c)
 \end{gathered}
\]
for all $a,b\in L$ and $c\in K$. Each condition defines a closed
set by continuity of evaluations and of the real arithmetic operations.
Conversely these conditions imply definiteness: if $a\ne0$, then
$f(a)f(a^{-1})=f(1)=1$, so $f(a)\ne0$.
They therefore define precisely an extending absolute value.
All these functions belong to $\prod_{a\in L}T_a$, which is compact
by Tychonoff's theorem. Their closed locus is compact.
\end{proof}

\begin{lemma}[Compact fibres and clopen cylinders]
\label{lem:compact-evaluation-fibre}
\leanok
\lean{ArakelovGeometry.AdelicStructure.Algebraic.instTopologicalSpaceBaseFiber,ArakelovGeometry.AdelicStructure.Algebraic.isEmbedding_fibreEval,ArakelovGeometry.AdelicStructure.Algebraic.instT2SpaceBaseFiber,ArakelovGeometry.AdelicStructure.Algebraic.instCompactSpaceBaseFiber,ArakelovGeometry.AdelicStructure.Algebraic.continuous_projFiniteFiber,ArakelovGeometry.AdelicStructure.Algebraic.isClopen_fibreCylinder,ArakelovGeometry.AdelicStructure.Algebraic.isCompact_fibreCylinder}
\source{chen-moriwaki-adelic:page-0208,chen-moriwaki-adelic:page-0209}
\uses{lem:compact-absolute-value-extensions,def:algebraic-projections,lem:finite-stage-discrete}
Give $X_\omega$ the topology induced by all evaluations
$x\mapsto |a|_x$ for $a\in L$. This topology is compact and Hausdorff.
Every finite-stage fibre has the discrete evaluation topology, every
restriction $p_{F,\omega}$ is continuous, and each cylinder
$p_{F,\omega}^{-1}(B)$ is clopen and compact.
\end{lemma}
\begin{proof}
The fixed base point identifies $X_\omega$ with the absolute values
extending $|\cdot|_\omega$. The evaluation map is injective, so its
induced topology is Hausdorff, and the preceding lemma makes it compact.
For a finite stage, evaluations again give an injective map into a
Hausdorff space. A finite Hausdorff space is discrete: each singleton
is the intersection of the complements of all other singletons.
Evaluating a restricted absolute value is evaluating the original
absolute value at the corresponding element of $L$. Thus all finite
restriction maps are continuous for the induced evaluation topologies.
Their target spaces are finite discrete, so inverse images of subsets
are clopen. Closed subsets of a compact space are compact.
\end{proof}

\begin{lemma}[Surjectivity on finite-stage fibres]\label{lem:finite-fibre-projection-surjective}
\leanok
\lean{ArakelovGeometry.AdelicStructure.Algebraic.surjective_projFiniteFiber}
\uses{prop:algebraic-restriction-surjective,lem:inseparable-extension-unique,def:algebraic-projections,def:finite-parameters}
\source{chen-moriwaki-adelic:page-0208}
Every projection $p_{F,\omega}:X_\omega\to X_{F,\omega}$ is surjective.
\end{lemma}
\begin{proof}
A point of $X_{F,\omega}$ gives an absolute value on the separable part
of $F/K$ extending $|\cdot|_\omega$. Extend it uniquely across the
purely inseparable extension to $F$, then extend that absolute value to
the algebraic extension $L/F$. The resulting value on $L$ still
restricts to $|\cdot|_\omega$ on $K$, and its projection is the given
point of the finite-stage fibre.
\end{proof}

\begin{lemma}[Compact inverse-limit fibres]\label{lem:compact-algebraic-fibre}\notready
\uses{def:algebraic-projections,lem:finite-stage-exhaustion,lem:finite-stage-discrete,lem:finite-stage-transitions,lem:finite-fibre-projection-surjective,lem:compact-evaluation-fibre}
\source{chen-moriwaki-adelic:page-0208,chen-moriwaki-adelic:page-0209}
Give each $X_{F,\omega}$ its finite discrete topology. Restriction
identifies $X_\omega$ with the closed subspace of
$\prod_{F\in\mathscr E_{L/K}}X_{F,\omega}$ consisting of compatible
families. In particular $X_\omega$ is compact, and every cylinder
$p_{F,\omega}^{-1}(B)$ is both open and compact.
\end{lemma}
\begin{proof}
Restrictions form compatible families, and the map to the product is
continuous because every coordinate projection is continuous. It is
injective: evaluating the coordinate at $K(a)$ recovers $|a|$ for every
$a\in L$, and these evaluations determine the absolute value.
For a compatible family $(x_F)$, put
$C_F=p_{F,\omega}^{-1}\{x_F\}$. Each $C_F$ is closed and nonempty,
by continuity and surjectivity of the finite projection. Compatibility
implies $C_{FG}\subseteq C_F\cap C_G$, so these sets have the finite
intersection property. The compactness established in
Lemma~\ref{lem:compact-evaluation-fibre} gives a point in their total
intersection, proving surjectivity onto compatible families.
Compatibility is the intersection, over inclusions $F\subseteq G$, of
the equations $x_F=x_G|_F$. Each equation is closed in the product of
finite discrete spaces. The resulting continuous bijection from the
compact space $X_\omega$ to this Hausdorff subspace is a homeomorphism.
A cylinder is the inverse image of a clopen subset of a finite discrete
space; it is clopen, and hence compact.
\end{proof}

\begin{remark}[Cylinder and Borel $\sigma$-algebras]\label{rem:cylinder-borel-distinction}
\uses{lem:compact-algebraic-fibre,def:projective-measurable}
\source{chen-moriwaki-adelic:page-0209}
The compactness argument uses the projective topology, while the measure
is constructed on the $\sigma$-algebra generated by cylinders. No
identification with the full Borel $\sigma$-algebra is required. For an
uncountable projective system the two $\sigma$-algebras need not coincide;
the assertion of their equality in the proof of Proposition 3.4.6 is
therefore omitted here.
\end{remark}

\begin{definition}[Cylinder functions]\label{def:algebraic-cylinders}
\uses{def:algebraic-projections}
\source{chen-moriwaki-adelic:page-0208,chen-moriwaki-adelic:page-0209}
\lean{ArakelovGeometry.AdelicStructure.Algebraic.DescendsToFinite,ArakelovGeometry.AdelicStructure.Algebraic.CylinderFunctions,ArakelovGeometry.AdelicStructure.Algebraic.DescendsToFiniteFiber,ArakelovGeometry.AdelicStructure.Algebraic.FibreCylinderFunctions}
A global real cylinder function is a function $f=g\circ p_F$ for some
finite stage $F$ and some function $g:\Omega_F\to\mathbb R$.
A fibre cylinder function is a function $f=g\circ p_{F,\omega}$ with
$g:X_{F,\omega}\to\mathbb R$. Write $V_L$ and $V_{L,\omega}$ for these
classes. A presentation consists of the stage $F$, the function $g$, and
the displayed factorization. The corresponding presentation spaces consist
of a function together with such a presentation; forgetting the presentation
has image $V_L$, respectively $V_{L,\omega}$.
\end{definition}

\begin{lemma}[Algebra of cylinder functions]\label{lem:cylinder-algebra}
\leanok
\uses{def:algebraic-cylinders,lem:finite-stage-exhaustion,lem:finite-stage-transitions}
\source{chen-moriwaki-adelic:page-0209}
\lean{ArakelovGeometry.AdelicStructure.Algebraic.DescendsToFinite.nonempty_add,ArakelovGeometry.AdelicStructure.Algebraic.DescendsToFinite.nonempty_smul,ArakelovGeometry.AdelicStructure.Algebraic.DescendsToFiniteFiber.nonempty_add,ArakelovGeometry.AdelicStructure.Algebraic.DescendsToFiniteFiber.nonempty_mul,ArakelovGeometry.AdelicStructure.Algebraic.DescendsToFiniteFiber.nonempty_smul}
Global cylinder functions are closed under addition and scalar
multiplication. Fibre cylinder functions are closed under addition,
multiplication, and scalar multiplication.
\end{lemma}
\begin{proof}
For presentations at $F$ and $G$, pull both functions back to $FG$.
Transitivity of restriction shows that their sum and product at $FG$
give the sum and product of the original functions. For scalar
multiplication multiply the descended function at the same stage.
\end{proof}

\begin{lemma}[Constant cylinder functions]\label{lem:cylinder-constants}\notready
\uses{def:algebraic-cylinders,def:finite-stages}
\source{chen-moriwaki-adelic:page-0209}
Every constant real function on $\Omega_L$ or on $X_\omega$ is a cylinder
function.
\end{lemma}
\begin{proof}
Represent the function by the same constant function at the stage $K$.
\end{proof}

\begin{definition}[Integral of a fibre cylinder]\label{def:cylinder-integral}
\leanok
\uses{def:algebraic-cylinders,def:finite-fibre-integral}
\source{chen-moriwaki-adelic:page-0208}
\lean{ArakelovGeometry.AdelicStructure.Algebraic.fibreCylinderIntegral}
For a presentation $f=g\circ p_{F,\omega}$ put
\[
 J_{F,g}(f)=\int_{X_{F,\omega}}g\,dP_{F,\omega}.
\]
Every such integral is finite, since the fibre is finite.
\end{definition}

\begin{lemma}[Independence of the cylinder presentation]\label{lem:cylinder-integral-independent}
\leanok
\uses{def:cylinder-integral,lem:finite-fibre-marginal,lem:finite-fibre-projection-surjective,lem:finite-stage-exhaustion}
\source{chen-moriwaki-adelic:page-0208,chen-moriwaki-adelic:page-0209}
\lean{ArakelovGeometry.AdelicStructure.Algebraic.fibreCylinderIntegral_congr}
If $f=g\circ p_{F,\omega}=h\circ p_{G,\omega}$, then
$J_{F,g}(f)=J_{G,h}(f)$.
\end{lemma}
\begin{proof}
Put $H=FG$. The two functions $g\circ p_{H/F}$ and $h\circ p_{H/G}$
on $X_{H,\omega}$ agree: lift any point to $X_\omega$ by surjectivity
and use the equality defining $f$. Finite marginal compatibility says
their integrals against $P_{H,\omega}$ are respectively $J_{F,g}$ and
$J_{G,h}$; hence these are equal.
\end{proof}

\begin{lemma}[The cylinder probability functional]\label{lem:cylinder-functional}\notready
\uses{lem:cylinder-integral-independent,lem:cylinder-algebra,lem:cylinder-constants,lem:algebraic-cylinder-integral,prop:algebraic-fibre-marginals}
\source{chen-moriwaki-adelic:page-0208,chen-moriwaki-adelic:page-0209}
The resulting functional $J_\omega$ on $V_{L,\omega}$ is positive,
linear, and sends $1$ to $1$.
\end{lemma}
\begin{proof}
A cylinder function is integrable on its fibre: the descended function
is integrable on the finite-stage probability space, and its pullback
has the same absolute integral by the marginal formula. By
Lemma~\ref{lem:algebraic-cylinder-integral}, its cylinder integral is
its integral against $P_{L,\omega}$. Positivity, additivity and scalar
linearity therefore follow from the corresponding properties of the
integral of integrable real functions. Every constant function has
integral equal to its value, because $P_{L,\omega}$ is a probability;
in particular $J_\omega(1)=1$. Independence of the presentation makes
these identities properties of the function itself.
\end{proof}

\begin{lemma}[Algebra of cylinder sets]\label{lem:fibre-cylinder-set-algebra}
\leanok
\uses{def:projective-measurable,lem:finite-stage-exhaustion,lem:finite-stage-transitions,lem:finite-stage-discrete}
\source{chen-moriwaki-adelic:page-0209}
\lean{ArakelovGeometry.AdelicStructure.Algebraic.fibreCylinders,ArakelovGeometry.AdelicStructure.Algebraic.isSetAlgebra_fibreCylinders,ArakelovGeometry.AdelicStructure.Algebraic.fibreMeasurableSpace_eq_generateFrom,ArakelovGeometry.AdelicStructure.Algebraic.measurable_projFiniteFiber}
The cylinders $\mathcal D_\omega=\{p_{F,\omega}^{-1}(B)\}$ form an
algebra of sets and generate $\mathcal A_{L,\omega}$.
\end{lemma}
\begin{proof}
The empty set is represented at the stage $K$. Complements are taken
at the same stage. To form a union, pass to the compositum of the two
stages and take the union of the two pulled-back subsets there.
Transitivity of the projections identifies its pullback with the
desired union. Since every subset of each finite-stage fibre is
measurable, making all finite-stage projections measurable is exactly
the requirement that every cylinder be measurable. This proves the
assertion about the generated $\sigma$-algebra.
\end{proof}

\begin{definition}[Cylinder content and outer measure]\label{def:cylinder-content}
\leanok
\uses{lem:fibre-cylinder-set-algebra,def:finite-fibre-integral}
\source{chen-moriwaki-adelic:page-0209}
\lean{ArakelovGeometry.AdelicStructure.Algebraic.fibreCylinderContent}
For any $A\subseteq X_\omega$, define
\[
 c_\omega(A)=\inf_{A\subseteq p_{F,\omega}^{-1}(B)}P_{F,\omega}(B),
 \qquad
 \mu_\omega^*(A)=\inf_{A\subseteq\bigcup_n A_n}\sum_n c_\omega(A_n).
\]
The first infimum runs over all finite stages and all subsets of their
fibres; the second runs over all countable covers by arbitrary subsets.
These formulas give a concrete realization of the Carath\'eodory
construction used in Proposition 3.4.6.
\end{definition}

\begin{lemma}[Cylinder premeasure]\label{lem:cylinder-premeasure}
\leanok
\uses{def:cylinder-content,lem:fibre-cylinder-set-algebra,lem:finite-fibre-marginal,lem:finite-fibre-projection-surjective,lem:compact-evaluation-fibre}
\source{chen-moriwaki-adelic:page-0209}
\lean{ArakelovGeometry.AdelicStructure.Algebraic.fibreCylinderContent_empty,ArakelovGeometry.AdelicStructure.Algebraic.finiteFibreProbabilityMeasure_le_of_preimage_subset,ArakelovGeometry.AdelicStructure.Algebraic.fibreCylinderContent_preimage,ArakelovGeometry.AdelicStructure.Algebraic.fibreCylinderContent_union,ArakelovGeometry.AdelicStructure.Algebraic.fibreCylinderAddContent,ArakelovGeometry.AdelicStructure.Algebraic.fibreCylinderAddContent_isSigmaSubadditive}
The restriction $p_\omega$ of $c_\omega$ to $\mathcal D_\omega$ is a
probability premeasure, with
$p_\omega(p_{F,\omega}^{-1}(B))=P_{F,\omega}(B)$.
For any sequence of cylinders whose union is a cylinder, one has
$p_\omega(\bigcup_n C_n)\leq\sum_n p_\omega(C_n)$.
\end{lemma}
\begin{proof}
Suppose $p_{F,\omega}^{-1}(B)\subseteq p_{G,\omega}^{-1}(D)$.
Pass to $H=FG$. Surjectivity of $p_{H,\omega}$ implies that the
pullback of $B$ in $X_{H,\omega}$ is contained in the pullback of $D$.
Monotonicity and finite marginal compatibility give
$P_{F,\omega}(B)\leq P_{G,\omega}(D)$. Taking the infimum over
cylinder supersets, and then using the cylinder itself as a superset,
proves $c_\omega(p_{F,\omega}^{-1}(B))=P_{F,\omega}(B)$.
In particular $c_\omega(\varnothing)=0$ and
$c_\omega(X_\omega)=1$.

For two disjoint cylinders, pass to a common stage. Surjectivity makes
their representing subsets disjoint; the additivity of the finite-stage
measure therefore proves finite additivity. If cylinders $C_n$ have
cylinder union $C$, they form an open cover of the compact set $C$.
A finite subcover and finite subadditivity give
$p_\omega(C)\leq\sum_n p_\omega(C_n)$.
When the $C_n$ are disjoint, every term outside a finite subcover is
empty. Finite additivity then gives the asserted countable additivity.
\end{proof}

\begin{lemma}[Outer measure of cylinders]\label{lem:cylinder-outer-measure}
\leanok
\uses{def:cylinder-content,lem:compact-evaluation-fibre,lem:cylinder-premeasure}
\source{chen-moriwaki-adelic:page-0209}
\lean{ArakelovGeometry.AdelicStructure.Algebraic.fibreCylinderOuterMeasure_eq_induced,ArakelovGeometry.AdelicStructure.Algebraic.fibreCylinderOuterMeasure_preimage,ArakelovGeometry.AdelicStructure.Algebraic.fibreMeasurableSpace_le_caratheodory}
One has $c_\omega(\varnothing)=0$. Moreover $\mu_\omega^*$ equals
the infimum of $\sum_n p_\omega(C_n)$ over countable cylinder covers,
every cylinder is Carath\'eodory measurable, and
$\mu_\omega^*(C)=p_\omega(C)$ for $C\in\mathcal D_\omega$.
\end{lemma}
\begin{proof}
The empty cylinder proves the first assertion. Let $\nu^*$ be the
infimum over countable cylinder covers. The countable-cover construction
is the largest outer measure bounded above by its assigned costs: if an
outer measure is bounded by each cost, its countable subadditivity
bounds it by every cover sum and hence by their infimum.
The preceding lemma gives $c_\omega(C)=p_\omega(C)$ on cylinders.
Since $\mu_\omega^*(C)\leq c_\omega(C)$, the maximality just proved
gives $\mu_\omega^*\leq\nu^*$. Conversely, for any $A\subseteq C$
with $C$ a cylinder, monotonicity gives
$\nu^*(A)\leq\nu^*(C)\leq p_\omega(C)$. Taking the infimum over
$C$ gives $\nu^*(A)\leq c_\omega(A)$, and maximality for the
arbitrary-set cover construction gives $\nu^*\leq\mu_\omega^*$.
Thus the two outer measures agree.

If cylinders $C_n$ cover a cylinder $C$, compactness gives a finite
subcover, and finite subadditivity gives
$p_\omega(C)\le\sum_n p_\omega(C_n)$. Covering $C$ by itself gives
the reverse inequality, so $\mu_\omega^*(C)=p_\omega(C)$.
Fix a cylinder $D$ and a cylinder cover $(C_n)$ of any $A$.
The cylinders $C_n\cap D$ and $C_n\setminus D$ cover respectively
$A\cap D$ and $A\setminus D$, and finite additivity yields
\[
 \mu_\omega^*(A\cap D)+\mu_\omega^*(A\setminus D)
 \le\sum_n p_\omega(C_n).
\]
Taking the infimum and using outer subadditivity proves the
Carath\'eodory equality. Thus every cylinder is measurable.
\end{proof}

\begin{definition}[Projective fibre probability]\label{def:algebraic-fibre-probability}
\leanok
\uses{lem:cylinder-outer-measure,def:projective-measurable}
\source{chen-moriwaki-adelic:page-0209}
\lean{ArakelovGeometry.AdelicStructure.Algebraic.fibreProbabilityMeasure}
Let $P_{L,\omega}$ be the restriction of $\mu_\omega^*$ to
$\mathcal A_{L,\omega}$. This restriction is a measure, since the
Carath\'eodory measurable sets form a $\sigma$-algebra containing the
cylinder generators.
\end{definition}

\begin{proposition}[Probability and finite marginals]\label{prop:algebraic-fibre-marginals}
\leanok
\uses{def:algebraic-fibre-probability,lem:cylinder-outer-measure}
\source{chen-moriwaki-adelic:page-0209}
\lean{ArakelovGeometry.AdelicStructure.Algebraic.fibreProbabilityMeasure_preimage,ArakelovGeometry.AdelicStructure.Algebraic.fibreProbabilityMeasure_isProbability,ArakelovGeometry.AdelicStructure.Algebraic.map_fibreProbabilityMeasure_projFiniteFiber}
$P_{L,\omega}(X_\omega)=1$ and
$(p_{F,\omega})_*P_{L,\omega}=P_{F,\omega}$ for every finite stage $F$.
\end{proposition}
\begin{proof}
The total space and all $p_{F,\omega}^{-1}(B)$ are cylinders. Their
outer measures are their prescribed premeasure values by
Lemma~\ref{lem:cylinder-outer-measure}, proving both assertions.
\end{proof}

\begin{lemma}[Uniqueness from finite marginals]\label{lem:algebraic-fibre-uniqueness}
\leanok
\uses{prop:algebraic-fibre-marginals,lem:cylinder-premeasure,def:projective-measurable}
\source{chen-moriwaki-adelic:page-0209}
\lean{ArakelovGeometry.AdelicStructure.Algebraic.fibreMeasure_ext,ArakelovGeometry.AdelicStructure.Algebraic.eq_fibreProbabilityMeasure_of_map_projFiniteFiber}
Let $\mu,\eta$ be measures on $(X_\omega,\mathcal A_{L,\omega})$,
with $\mu$ finite. If $(p_{F,\omega})_*\mu=(p_{F,\omega})_*\eta$
for every finite stage $F$, then $\mu=\eta$.
In particular any measure whose finite marginals are $P_{F,\omega}$
equals $P_{L,\omega}$; its total mass is automatically one.
\end{lemma}
\begin{proof}
For a cylinder $A=p_{F,\omega}^{-1}(B)$, measurability of the finite-stage
projection gives
$\mu(A)=((p_{F,\omega})_*\mu)(B)=((p_{F,\omega})_*\eta)(B)=\eta(A)$.
The total space is a cylinder, so the measures have the same finite
total mass. The sets on which they agree form a Dynkin system:
complements follow by subtracting from that common finite mass, and
disjoint countable unions follow by countable additivity. Cylinders
form an algebra and generate $\mathcal A_{L,\omega}$. The
$\pi$-$\lambda$ theorem therefore proves $\mu=\eta$.
Apply this finite-measure uniqueness argument with one measure equal
to $P_{L,\omega}$ and use its prescribed marginals. Equality also gives
total mass one for the other measure.
\end{proof}

\begin{lemma}[Integrals of fibre cylinders]\label{lem:algebraic-cylinder-integral}
\leanok
\uses{prop:algebraic-fibre-marginals,def:cylinder-integral,lem:finite-stage-discrete}
\source{chen-moriwaki-adelic:page-0209}
\lean{ArakelovGeometry.AdelicStructure.Algebraic.integral_fibreProbabilityMeasure_descend}
For a fibre cylinder $f=g\circ p_{F,\omega}$,
$\int f\,dP_{L,\omega}=J_\omega(f)$.
\end{lemma}
\begin{proof}
Write $g$ as the finite sum of its values times singleton indicators.
Pullback gives the corresponding sum of cylinder indicators.
Integrating and applying the finite marginal formula gives exactly
$\sum_{y\in X_{F,\omega}}g(y)P_{F,\omega}(\{y\})$ on both sides.
\end{proof}

\begin{definition}[Algebraic fibre integration]\label{def:algebraic-fibre-integral}
\leanok
\uses{def:algebraic-fibre-probability}
\source{chen-moriwaki-adelic:page-0209,chen-moriwaki-adelic:page-0210}
\lean{ArakelovGeometry.AdelicStructure.Algebraic.fibreLIntegral,ArakelovGeometry.AdelicStructure.Algebraic.fibreIntegral}
For $f:\Omega_L\to[0,\infty]$ set
$I^+_{L/K}f(\omega)=\int_{X_\omega}f\,dP_{L,\omega}$.
For real $f$, set $I_{L/K}f(\omega)$ to its real integral on the fibre
when integrable, and to $0$ otherwise. This last convention makes a
total function; identities for signed integrals below specify the
hypotheses under which it agrees with ordinary integration.
\end{definition}

\begin{lemma}[Integration of pullbacks]\label{lem:algebraic-integral-pullback}
\leanok
\uses{def:algebraic-fibre-integral,lem:algebraic-cylinder-integral,prop:algebraic-fibre-marginals}
\source{chen-moriwaki-adelic:page-0209,chen-moriwaki-adelic:page-0211}
\lean{ArakelovGeometry.AdelicStructure.Algebraic.fibreIntegral_comp_proj,ArakelovGeometry.AdelicStructure.Algebraic.fibreIntegral_comp_projFinite}
For every real $g:\Omega\to\mathbb R$,
$I_{L/K}(g\circ\pi_{L/K})=g$. For every finite stage $F$ and every
real $h:\Omega_F\to\mathbb R$,
$I_{L/K}(h\circ p_F)=I_{F/K}h$.
\end{lemma}
\begin{proof}
The first integrand is the constant $g(\omega)$ on a probability fibre.
For the second apply Lemma~\ref{lem:algebraic-cylinder-integral} on
each fibre. No global measurability is needed, since every real function
on a finite fibre is measurable and integrable.
\end{proof}

\begin{lemma}[Nonnegative integration of pullbacks]\label{lem:algebraic-lintegral-pullback}
\leanok
\uses{def:algebraic-fibre-integral,prop:algebraic-fibre-marginals,lem:finite-stage-discrete}
\source{chen-moriwaki-adelic:page-0209,chen-moriwaki-adelic:page-0211}
\lean{ArakelovGeometry.AdelicStructure.Algebraic.fibreLIntegral_comp_proj,ArakelovGeometry.AdelicStructure.Algebraic.fibreLIntegral_comp_projFinite}
For every $g:\Omega\to[0,\infty]$, $I^+_{L/K}(g\circ\pi_{L/K})=g$.
For every finite stage $F$ and every $h:\Omega_F\to[0,\infty]$,
$I^+_{L/K}(h\circ p_F)=I^+_{F/K}h$.
\end{lemma}
\begin{proof}
The first integrand is constant on a probability fibre.
For the second the marginal formula gives
equality first for indicators, then finite nonnegative sums, and then
all such functions by increasing simple approximation and monotone
convergence. No global measurability of $g$ or $h$ is needed here,
since every restriction to a finite fibre is measurable.
\end{proof}

\begin{lemma}[Restriction of measurable functions to fibres]\label{lem:algebraic-fibre-measurable}
\leanok
\uses{def:projective-measurable,def:algebraic-projections}
\source{chen-moriwaki-adelic:page-0210}
\lean{ArakelovGeometry.AdelicStructure.Algebraic.measurable_subtype_val_baseFiber,ArakelovGeometry.AdelicStructure.Algebraic.measurable_restrict_baseFiber}
The inclusion $X_\omega\to\Omega_L$ is measurable. In particular any
$\mathcal A_L$-measurable function with values in a measurable space
restricts to an $\mathcal A_{L,\omega}$-measurable function.
\end{lemma}
\begin{proof}
The inverse image of a generating set $p_F^{-1}(B)$ is
$p_{F,\omega}^{-1}(B\cap X_{F,\omega})$. The intersection is measurable
for the induced measurable structure on the fibre, and the fibre
projection is measurable by definition. Closure of inverse
images under complements and countable unions proves the assertion
for the generated $\sigma$-algebra. Composition proves the assertion
for functions.
\end{proof}

\begin{lemma}[Measurable finite-stage transitions]\label{lem:finite-stage-transitions-measurable}
\leanok
\uses{lem:finite-stage-transitions,lem:finite-measurable-space}
\source{chen-moriwaki-adelic:page-0208,chen-moriwaki-adelic:page-0210}
\lean{ArakelovGeometry.AdelicStructure.Algebraic.measurable_projFiniteOfLE}
For finite intermediate fields $F\subseteq G$, the transition
$p_{G/F}:\Omega_G\to\Omega_F$ is measurable.
\end{lemma}
\begin{proof}
The measurable structure on $\Omega_F$ is generated by its base
coordinate and the absolute-value evaluations of elements of
$F_{\mathrm{sep}}$. The transition preserves the base coordinate.
For $a\in F_{\mathrm{sep}}$, its evaluation after transition is the
evaluation of its image in $G_{\mathrm{sep}}$. Both types of functions
are measurable on $\Omega_G$, proving the assertion.
\end{proof}

\begin{lemma}[Measurable family of fibre probabilities]\label{lem:algebraic-measurable-probability-family}
\leanok
\uses{lem:algebraic-fibre-measurable,lem:finite-stage-transitions-measurable,prop:algebraic-fibre-marginals,lem:finite-kernel-measurable,lem:finite-stage-exhaustion,lem:projective-evaluations}
\source{chen-moriwaki-adelic:page-0210}
\lean{ArakelovGeometry.AdelicStructure.Algebraic.measurable_map_fibreProbabilityMeasure}
Let $i_\omega:X_\omega\to\Omega_L$ denote inclusion and set
$\kappa_\omega=(i_\omega)_*P_{L,\omega}$. Then $\kappa_\omega$
is a probability measure on $(\Omega_L,\mathcal A_L)$, and for every
$A\in\mathcal A_L$, the function $\omega\mapsto\kappa_\omega(A)$
is measurable.
\end{lemma}
\begin{proof}
The inclusion is measurable and preserves total mass under pushforward,
so $\kappa_\omega$ is a probability. For a measurable finite-stage
cylinder $A=p_F^{-1}(B)$, commutation of the projections and the finite
marginal formula give
\[
 \kappa_\omega(A)=P_{F,\omega}(B\cap X_{F,\omega})
 =I^+_{F/K}(\mathbf1_B)(\omega),
\]
which is measurable by finite-kernel measurability. These cylinders
form a $\pi$-system: two finite stages embed in their finite compositum,
and measurable transitions express their intersection at that common
stage. They generate $\mathcal A_L$ by its definition.
The measurable sets $A$ for which $\omega\mapsto\kappa_\omega(A)$
is measurable form a Dynkin system. The total space has value $1$;
the complement of $A$ has value $1-\kappa_\omega(A)$; a disjoint
countable union has the pointwise sum of the corresponding values.
These operations preserve measurability, since all individual
probabilities are finite. The $\pi$-$\lambda$ theorem completes the proof.
\end{proof}

\begin{proposition}[Measurability of fibre integration]\label{prop:algebraic-kernel-measurable}
\leanok
\uses{lem:algebraic-fibre-measurable,lem:algebraic-measurable-probability-family,def:algebraic-fibre-integral}
\source{chen-moriwaki-adelic:page-0210}
\lean{ArakelovGeometry.AdelicStructure.Algebraic.measurable_fibreLIntegral}
If $f:\Omega_L\to[0,\infty]$ is $\mathcal A_L$-measurable, then
$I^+_{L/K}f$ is $\mathcal A$-measurable.
\end{proposition}
\begin{proof}
Write $\kappa_\omega=(i_\omega)_*P_{L,\omega}$ as in
Lemma~\ref{lem:algebraic-measurable-probability-family}. Integrals of
nonnegative measurable simple functions against $\kappa_\omega$
are measurable functions of $\omega$, by finite sums of measurable
evaluations $\kappa_\omega(A)$. Increasing simple approximation and
monotone convergence give measurability of $\omega\mapsto\int f\,d\kappa_\omega$.
The pushforward integral formula identifies this with
$\int_{X_\omega}f\circ i_\omega\,dP_{L,\omega}=I^+_{L/K}f(\omega)$.
\end{proof}

\begin{definition}[Measure on the algebraic extension]\label{def:algebraic-measure}
\leanok
\uses{prop:algebraic-kernel-measurable}
\source{chen-moriwaki-adelic:page-0210}
\lean{ArakelovGeometry.AdelicStructure.Algebraic.measure,ArakelovGeometry.AdelicStructure.Algebraic.measure_apply}
For $B\in\mathcal A_L$ define
\[
 \nu_L(B)=\int_\Omega I^+_{L/K}(\mathbf1_B)(\omega)\,d\nu(\omega).
\]
This prescription is a measure on $(\Omega_L,\mathcal A_L)$.
\end{definition}
\begin{proof}
The empty set has zero fibre integral. For disjoint measurable $B_n$,
their intersections with each fibre are disjoint and measurable, so
the fibre probability of their union is the sum of their fibre
probabilities. These summands are measurable by
Proposition~\ref{prop:algebraic-kernel-measurable}. Monotone convergence
of the partial sums under $\nu$ proves countable additivity.
\end{proof}

\begin{proposition}[Nonnegative disintegration]\label{prop:algebraic-lintegral}
\leanok
\uses{def:algebraic-measure,prop:algebraic-kernel-measurable}
\source{chen-moriwaki-adelic:page-0210}
\lean{ArakelovGeometry.AdelicStructure.Algebraic.lintegral_measure}
For measurable $f:\Omega_L\to[0,\infty]$,
\[
 \int_{\Omega_L} f\,d\nu_L=\int_\Omega I^+_{L/K}f\,d\nu.
\]
\end{proposition}
\begin{proof}
The assertion for an indicator is the definition of $\nu_L$.
Additivity and multiplication by a nonnegative scalar prove it for
nonnegative simple functions. Choose simple $f_n\uparrow f$.
Monotone convergence applies under $\nu_L$, under every fibre
probability, and then under $\nu$, proving the equality, including
the value $+\infty$.
\end{proof}

\begin{proposition}[Integrability criterion]\label{prop:algebraic-integrability}
\leanok
\uses{prop:algebraic-lintegral}
\source{chen-moriwaki-adelic:page-0210,chen-moriwaki-adelic:page-0211}
\lean{ArakelovGeometry.AdelicStructure.Algebraic.integrable_measure_iff}
A measurable real function $f$ on $\Omega_L$ is $\nu_L$-integrable
if and only if
\[
 \int_\Omega\left(\int_{X_\omega}|f|\,dP_{L,\omega}\right)d\nu<\infty,
\]
where both integrals in this criterion take values in $[0,\infty]$.
\end{proposition}
\begin{proof}
For measurable real $f$, integrability is precisely finiteness of its
nonnegative absolute-value integral. Apply
Proposition~\ref{prop:algebraic-lintegral} to $|f|$. In particular,
finiteness is asserted before converting any integral to a real number.
\end{proof}

\begin{lemma}[Almost everywhere statements on algebraic fibres]\label{lem:algebraic-ae-fibres}
\leanok
\source{chen-moriwaki-adelic:page-0210,chen-moriwaki-adelic:page-0211}
\uses{def:algebraic-measure,prop:algebraic-kernel-measurable,lem:algebraic-fibre-measurable}
\lean{ArakelovGeometry.AdelicStructure.Algebraic.ae_ae_fibre_of_ae,ArakelovGeometry.AdelicStructure.Algebraic.ae_fibre_measure_preimage_eq_zero}
If a property holds for $\nu_L$-almost every point, it holds
$P_{L,\omega}$-almost everywhere on $X_\omega$ for $\nu$-almost every
$\omega$. In particular, every $\nu_L$-null set has null intersection
with almost every fibre.
\end{lemma}
\begin{proof}
Enclose the exceptional set in a measurable null set $N$.
The defining measure formula gives
$\int_\Omega P_{L,\omega}(N\cap X_\omega)\,d\nu(\omega)=0$.
The integrand is nonnegative and measurable, so it vanishes almost
everywhere. On each such fibre, the exceptional subset of $N$ is null
by monotonicity. Applying this to the property of lying outside a
given null set proves the last assertion.
\end{proof}

\begin{lemma}[Almost everywhere integrability on fibres]\label{lem:algebraic-measurable-fibre-integrable}
\leanok
\source{chen-moriwaki-adelic:page-0210,chen-moriwaki-adelic:page-0211}
\uses{prop:algebraic-integrability,prop:algebraic-kernel-measurable,lem:algebraic-fibre-measurable}
\lean{ArakelovGeometry.AdelicStructure.Algebraic.integrable_fibre_ae_of_measurable}
If $f$ is measurable and $\nu_L$-integrable, its restriction to
$X_\omega$ is $P_{L,\omega}$-integrable for $\nu$-almost every $\omega$.
\end{lemma}
\begin{proof}
The criterion gives $\int_\Omega I^+_{L/K}|f|\,d\nu<\infty$.
Thus $I^+_{L/K}|f|(\omega)<\infty$ almost everywhere. The restriction
of $f$ is measurable on every fibre; at the indicated base points
its absolute-value integral is finite, which is precisely integrability.
\end{proof}

\begin{lemma}[Nonnegative real fibre averages]\label{lem:algebraic-nonnegative-real-average}
\leanok
\source{chen-moriwaki-adelic:page-0210,chen-moriwaki-adelic:page-0211}
\uses{prop:algebraic-integrability,prop:algebraic-lintegral,prop:algebraic-kernel-measurable,lem:algebraic-fibre-measurable}
\lean{ArakelovGeometry.AdelicStructure.Algebraic.fibreIntegral_eq_toReal_of_nonneg,ArakelovGeometry.AdelicStructure.Algebraic.integrable_fibreIntegral_of_nonneg,ArakelovGeometry.AdelicStructure.Algebraic.integral_measure_of_nonneg}
For nonnegative measurable real $f$, its real fibre integral is the
real value of $I^+_{L/K}f$ whenever that value is finite. If the
nonnegative integral is infinite, the totalized real integral is zero.
If in addition $f$ is $\nu_L$-integrable, then $I_{L/K}f$ is integrable
on $\Omega$ and
$\int_{\Omega_L}f\,d\nu_L=\int_\Omega I_{L/K}f\,d\nu$.
\end{lemma}
\begin{proof}
On a fibre with finite nonnegative integral, the ordinary real integral
equals that value. On a fibre with infinite integral, the nonnegative
function is not integrable and its totalized real integral is zero.
If $f$ is globally integrable, nonnegative disintegration shows that
$I^+_{L/K}f$ has finite integral on the base. It is therefore finite
almost everywhere, and its real value is integrable with the same
integral. Discarding the null set where it is infinite proves the
stated real integration formula.
\end{proof}

\begin{lemma}[Signed integration for measurable functions]\label{lem:algebraic-measurable-signed-disintegration}
\leanok
\source{chen-moriwaki-adelic:page-0211}
\uses{lem:algebraic-nonnegative-real-average,lem:algebraic-measurable-fibre-integrable}
\lean{ArakelovGeometry.AdelicStructure.Algebraic.integrable_fibreIntegral_and_integral_of_measurable}
If $f$ is measurable and $\nu_L$-integrable, then $I_{L/K}f$ is
$\nu$-integrable and
$\int_{\Omega_L}f\,d\nu_L=\int_\Omega I_{L/K}f\,d\nu$.
\end{lemma}
\begin{proof}
Write $f=f^+-f^-$, where both parts are nonnegative, measurable and
integrable. The preceding lemma applies to each part. Their fibre
restrictions are integrable outside the union of two null sets.
There integral subtraction gives $I_{L/K}f=I_{L/K}f^+-I_{L/K}f^-$.
The difference is integrable on the base. Integrating this equality
and subtracting the two nonnegative integration formulas proves the result.
\end{proof}

\begin{lemma}[Null sets and measurable representatives]\label{lem:algebraic-null-disintegration}
\leanok
\uses{prop:algebraic-lintegral,prop:algebraic-kernel-measurable}
\source{chen-moriwaki-adelic:page-0211}
\lean{ArakelovGeometry.AdelicStructure.Algebraic.ae_fibre_measure_preimage_eq_zero,ArakelovGeometry.AdelicStructure.Algebraic.exists_measurable_representative,ArakelovGeometry.AdelicStructure.Algebraic.ae_ae_fibre_of_ae}
If a measurable $N\subseteq\Omega_L$ is $\nu_L$-null, then
$P_{L,\omega}(N\cap X_\omega)=0$ for $\nu$-almost every $\omega$.
Consequently, if $f$ is $\nu_L$-integrable, it admits a measurable
representative $\widetilde f=f$ almost everywhere such that, for
almost every $\omega$, $f$ and $\widetilde f$ agree
$P_{L,\omega}$-almost everywhere on $X_\omega$.
\end{lemma}
\begin{proof}
Disintegration of $\mathbf1_N$ gives a nonnegative measurable
function on the base with integral zero; hence that function vanishes
almost everywhere. An integrable real function is almost everywhere
equal to a measurable function. The exceptional set is contained in
a measurable null set by the definition of almost everywhere equality.
Apply the first assertion to that null set. Its fibre intersections
contain the fibre exceptional sets, proving the claim.
\end{proof}

\begin{lemma}[Fibre integration respects almost everywhere equality]
\label{lem:algebraic-fibre-integral-ae}
\leanok
\uses{lem:algebraic-null-disintegration}
\source{chen-moriwaki-adelic:page-0211}
\lean{ArakelovGeometry.AdelicStructure.Algebraic.fibreIntegral_congr_ae}
If real functions $f,g$ on $\Omega_L$ are equal $\nu_L$-almost everywhere,
then $I_{L/K}f=I_{L/K}g$ $\nu$-almost everywhere.
\end{lemma}
\begin{proof}
Their restrictions are equal $P_{L,\omega}$-almost everywhere for
$\nu$-almost every $\omega$ by
Lemma~\ref{lem:algebraic-null-disintegration}. On each such fibre their
real integrals agree, since the integral is unchanged by alteration
on a null set. This also applies when these functions are not integrable,
as both integrals then use the same totalized convention.
\end{proof}

\begin{lemma}[Integrability of signed fibre averages]\label{lem:algebraic-signed-integrability}
\leanok
\uses{prop:algebraic-integrability,prop:algebraic-kernel-measurable,lem:algebraic-ae-fibres,lem:algebraic-measurable-signed-disintegration}
\source{chen-moriwaki-adelic:page-0211}
\lean{ArakelovGeometry.AdelicStructure.Algebraic.integrable_fibre_ae,ArakelovGeometry.AdelicStructure.Algebraic.integrable_fibreIntegral,ArakelovGeometry.AdelicStructure.Algebraic.ofReal_abs_fibreIntegral_le_fibreLIntegral_ae}
For $\nu_L$-integrable real $f$, its fibre restriction is integrable
for $\nu$-almost every $\omega$, the function $I_{L/K}f$ is
$\nu$-integrable, and $|I_{L/K}f|\le I^+_{L/K}|f|$ almost everywhere.
\end{lemma}
\begin{proof}
First suppose $f$ is measurable. The measurable nonnegative function
$A=I^+_{L/K}|f|$ has finite integral, so $A<\infty$ almost everywhere.
On that set, the fibre is integrable and
$I_{L/K}f=I^+_{L/K}f^+-I^+_{L/K}f^-$, while
$|I_{L/K}f|\le A$. Define the displayed difference where $A<\infty$
and zero elsewhere. It is a measurable representative of $I_{L/K}f$,
and domination by $A$ proves integrability and the asserted bound.

For general integrable $f$, choose a measurable representative equal
to $f$ almost everywhere. Lemma~\ref{lem:algebraic-ae-fibres} makes
their restrictions equal almost everywhere on almost every fibre.
Their fibre integrals, global integrals and absolute integrals therefore
agree almost everywhere, so all assertions transfer.
\end{proof}

\begin{proposition}[Signed disintegration and contraction bound]\label{prop:algebraic-signed-disintegration}
\leanok
\uses{lem:algebraic-signed-integrability,prop:algebraic-lintegral,lem:algebraic-ae-fibres,lem:algebraic-measurable-signed-disintegration}
\source{chen-moriwaki-adelic:page-0211}
\lean{ArakelovGeometry.AdelicStructure.Algebraic.integral_abs_fibreIntegral_le,ArakelovGeometry.AdelicStructure.Algebraic.integral_measure}
For $\nu_L$-integrable real $f$,
\[
 \int_\Omega|I_{L/K}f|\,d\nu\le\int_{\Omega_L}|f|\,d\nu_L,
 \qquad
 \int_{\Omega_L}f\,d\nu_L=\int_\Omega I_{L/K}f\,d\nu.
\]
\end{proposition}
\begin{proof}
For measurable $f$, integrate the bound in
Lemma~\ref{lem:algebraic-signed-integrability} and apply nonnegative
disintegration to $|f|$. Apply the same disintegration to $f^+$ and
$f^-$, whose integrals are finite, and subtract to obtain the second
identity. For general integrable $f$, use its measurable representative
and Lemma~\ref{lem:algebraic-ae-fibres}; both global and fibre
integrals are unchanged almost everywhere.
\end{proof}

\begin{corollary}[The operator on $L^1$]\label{cor:algebraic-L1-operator}
\leanok
\uses{prop:algebraic-signed-disintegration,lem:algebraic-signed-integrability,lem:algebraic-fibre-integral-ae}
\source{chen-moriwaki-adelic:page-0211}
\lean{ArakelovGeometry.AdelicStructure.Algebraic.fibreIntegral_add_ae,ArakelovGeometry.AdelicStructure.Algebraic.fibreIntegral_smul,ArakelovGeometry.AdelicStructure.Algebraic.fibreIntegralL1,ArakelovGeometry.AdelicStructure.Algebraic.coeFn_fibreIntegralL1,ArakelovGeometry.AdelicStructure.Algebraic.fibreIntegralL1_add,ArakelovGeometry.AdelicStructure.Algebraic.fibreIntegralL1_smul,ArakelovGeometry.AdelicStructure.Algebraic.norm_fibreIntegralL1_le,ArakelovGeometry.AdelicStructure.Algebraic.fibreIntegralL1ContinuousLinearMap,ArakelovGeometry.AdelicStructure.Algebraic.norm_fibreIntegralL1ContinuousLinearMap_le,ArakelovGeometry.AdelicStructure.Algebraic.fibreIntegralL1_toL1}
Fibre integration induces a linear contraction from
$L^1(\Omega_L,\nu_L)$ to $L^1(\Omega,\nu)$.
Precisely, there is a continuous linear map $T$ satisfying
$T[f]=[I_{L/K}f]$ for every integrable real $f$, and $\|T\|\le1$.
\end{corollary}
\begin{proof}
By Lemma~\ref{lem:algebraic-signed-integrability}, the fibre average of
an integrable function is integrable. Define $T[f]=[I_{L/K}f]$.
Lemma~\ref{lem:algebraic-fibre-integral-ae} proves that this class is
independent of the chosen representative of $[f]$.

If $f,g$ are integrable, their restrictions to each fibre are integrable
outside the union of two null sets of base points. On all other fibres,
integral additivity gives $I_{L/K}(f+g)=I_{L/K}f+I_{L/K}g$.
For every scalar $c$, homogeneity of the real integral gives
$I_{L/K}(cf)=cI_{L/K}f$; this identity is valid also on exceptional
fibres under the totalized integral convention. These identities prove
linearity of $T$ on the equivalence classes.

Finally, Proposition~\ref{prop:algebraic-signed-disintegration} gives
\[
 \|T[f]\|_1=\int_\Omega|I_{L/K}f|\,d\nu
 \le\int_{\Omega_L}|f|\,d\nu_L=\|[f]\|_1.
\]
Thus $T$ is continuous and has operator norm at most one, with the
asserted formula for every integrable representative.
\end{proof}

\begin{lemma}[Finite-stage pushforward]\label{lem:algebraic-measure-marginals}
\leanok
\uses{def:algebraic-measure,prop:algebraic-fibre-marginals,thm:finite-lintegral,lem:projective-evaluations}
\source{chen-moriwaki-adelic:page-0211}
\lean{ArakelovGeometry.AdelicStructure.Algebraic.map_measure_projFinite}
For every finite stage $F$, $(p_F)_*\nu_L=\nu_F$.
\end{lemma}
\begin{proof}
For measurable $B\subseteq\Omega_F$, the defining formula and the
fibre marginal give
\[
 \nu_L(p_F^{-1}B)=\int_\Omega P_{F,\omega}(B\cap X_{F,\omega})\,d\nu
 =\nu_F(B).
\]
The last equality is finite disintegration of $\mathbf1_B$.
Measurability of $p_F$ identifies the left side with its pushforward.
\end{proof}

\begin{lemma}[Finite projections along field embeddings]
\label{lem:algebraic-finite-embedding-projection}
\leanok
\source{chen-moriwaki-adelic:page-0209,chen-moriwaki-adelic:page-0211}
\uses{prop:algebraic-fibre-marginals,lem:finite-fibre-embedding-naturality,lem:projective-evaluations,prop:projective-measurable-equality,def:algebraic-measure,prop:algebraic-lintegral,thm:finite-lintegral}
\lean{ArakelovGeometry.AdelicStructure.Algebraic.projFiniteAlong,ArakelovGeometry.AdelicStructure.Algebraic.projFiniteAlong_base,ArakelovGeometry.AdelicStructure.Algebraic.projFiniteAlong_ϕ,ArakelovGeometry.AdelicStructure.Algebraic.measurable_projFiniteAlong,ArakelovGeometry.AdelicStructure.Algebraic.projFiniteFiberAlong,ArakelovGeometry.AdelicStructure.Algebraic.projFiniteFiberAlong_eq_homFiber,ArakelovGeometry.AdelicStructure.Algebraic.measurable_projFiniteFiberAlong,ArakelovGeometry.AdelicStructure.Algebraic.map_fibreProbabilityMeasure_projFiniteFiberAlong,ArakelovGeometry.AdelicStructure.Algebraic.lintegral_projFiniteFiberAlong,ArakelovGeometry.AdelicStructure.Algebraic.fibreLIntegral_comp_projFiniteAlong,ArakelovGeometry.AdelicStructure.Algebraic.map_measure_projFiniteAlong}
Let $L/K$ be algebraic, let $E/K$ be finite, and let
$i:E\hookrightarrow L$ be a $K$-embedding. Pullback of absolute values
defines a measurable map $p_i:\Omega_L\to\Omega_E$, preserving the
base point, with $|a|_{p_i(w)}=|i(a)|_w$. Its restriction $p_{i,\omega}$
to each fibre is measurable and satisfies
\[
 (p_{i,\omega})_*P_{L,\omega}=P_{E,\omega},\qquad
 (p_i)_*\nu_L=\nu_E.
\]
For every function $g:\Omega_E\to[0,\infty]$, without a measurability
requirement away from a fibre,
$I^+_{L/K}(g\circ p_i)=I^+_{E/K}g$.
\end{lemma}
\begin{proof}
Put $F=i(E)$. This is a finite intermediate field of $L/K$, and $i$
identifies $E$ with $F$ over $K$. The map $p_i$ is the finite-stage
projection to $F$ followed by pullback along this isomorphism.
Equivalently, it is characterized by its base point and the displayed
evaluation identity. These coordinates are measurable, proving the
global measurability assertion. On each fibre it factors as the
measurable projection to the finite $F$-fibre followed by a map between
finite discrete fibres. The prescribed marginal at $F$ and
Lemma~\ref{lem:finite-fibre-embedding-naturality} give its probability
pushforward. Every function on the finite target fibre is measurable,
so change of variables proves the nonnegative fibre-integral identity.
Finally, apply that identity to indicators of measurable subsets of
$\Omega_E$ and integrate over the base. The defining disintegration
formulas for $\nu_L$ and $\nu_E$ give the global pushforward equality.
\end{proof}

\begin{lemma}[Base pushforward]\label{lem:algebraic-base-pushforward}
\leanok
\uses{def:algebraic-measure,prop:algebraic-fibre-marginals,lem:projective-evaluations}
\source{chen-moriwaki-adelic:page-0211}
\lean{ArakelovGeometry.AdelicStructure.Algebraic.map_measure_proj}
One has $(\pi_{L/K})_*\nu_L=\nu$.
\end{lemma}
\begin{proof}
For measurable $B\subseteq\Omega$, the fibre probability of
$\pi^{-1}B$ is $\mathbf1_B(\omega)$. The definition of $\nu_L$
therefore gives $\nu_L(\pi^{-1}B)=\nu(B)$.
\end{proof}

\begin{proposition}[Pullback of integrable functions]\label{prop:algebraic-base-pullback}
\leanok
\uses{lem:algebraic-base-pushforward,prop:algebraic-signed-disintegration,lem:algebraic-integral-pullback}
\source{chen-moriwaki-adelic:page-0211}
\lean{ArakelovGeometry.AdelicStructure.Algebraic.integrable_comp_proj_measure,ArakelovGeometry.AdelicStructure.Algebraic.integral_comp_proj_measure}
If $g\in L^1(\Omega,\nu)$, then $g\circ\pi_{L/K}$ is $\nu_L$-integrable
and its integral is $\int_\Omega g\,d\nu$.
\end{proposition}
\begin{proof}
The pushforward identities give equality of the nonnegative absolute
integrals for measurable representatives. Pullbacks of null
exceptional sets are null by the same identities, so pullback is
almost everywhere measurable and integrable for general integrable
functions. Applying the pushforward integration formula to the positive
and negative parts proves equality of signed integrals.
\end{proof}

\begin{lemma}[Finite-stage pullback on $L^1$]\label{lem:algebraic-finite-pullback}
\leanok
\uses{lem:algebraic-measure-marginals}
\source{chen-moriwaki-adelic:page-0211}
\lean{ArakelovGeometry.AdelicStructure.Algebraic.integrable_comp_projFinite_measure,ArakelovGeometry.AdelicStructure.Algebraic.integral_comp_projFinite_measure,ArakelovGeometry.AdelicStructure.Algebraic.integral_norm_comp_projFinite_measure}
For a finite stage $F$ and a real $\nu_F$-integrable function $g$,
the pullback $g\circ p_F$ is $\nu_L$-integrable, its integral equals
$\int g\,d\nu_F$, and its $L^1$ norm equals $\int |g|\,d\nu_F$.
\end{lemma}
\begin{proof}
The identity $(p_F)_*\nu_L=\nu_F$ gives equality of the nonnegative
integrals of $|g|$, $g^+$, and $g^-$ for measurable $g$.
It also pulls measurable null sets back to null sets. Thus a measurable
representative of an integrable function pulls back to a measurable
representative with the same finite absolute integral. Subtracting the
positive and negative integral equalities proves the signed assertion.
\end{proof}

\begin{proposition}[Integrable logarithmic evaluations]\label{prop:algebraic-log-integrable}
\leanok
\uses{lem:finite-stage-exhaustion,lem:finite-log-integrable,lem:algebraic-finite-pullback}
\source{chen-moriwaki-adelic:page-0211}
\lean{ArakelovGeometry.AdelicStructure.Algebraic.integrable_log_measure}
For $a\in L^\times$, the function $x\mapsto\log|a|_x$ is
$\nu_L$-integrable.
\end{proposition}
\begin{proof}
Take $F=K(a)$. The finite-stage logarithmic evaluation is integrable
by Lemma~\ref{lem:finite-log-integrable}. Its pullback along $p_F$
is the logarithmic evaluation on $L$, and
Lemma~\ref{lem:algebraic-finite-pullback} proves its integrability.
\end{proof}

\begin{definition}[Algebraic extension of an adelic structure]\label{def:algebraic-extension}
\leanok
\uses{def:adelic-structure,def:algebraic-measure,lem:projective-evaluations,prop:algebraic-log-integrable}
\source{chen-moriwaki-adelic:page-0212}
\lean{ArakelovGeometry.AdelicStructure.algebraicExtension}
The adelic extension $S\otimes_K L$ is
$(L,(\Omega_L,\mathcal A_L,\nu_L),\phi_L)$, where $\phi_L(\omega,w)=w$.
The measurable and integrable logarithmic evaluations just proved
make this an adelic structure.
\end{definition}

\begin{proposition}[Preservation of properness]\label{prop:algebraic-proper}
\leanok
\uses{def:algebraic-extension,prop:finite-proper,lem:algebraic-finite-pullback}
\source{chen-moriwaki-adelic:page-0211}
\lean{ArakelovGeometry.AdelicStructure.Algebraic.isProper_algebraicExtension}
If $S$ is proper, then $S\otimes_K L$ is proper.
\end{proposition}
\begin{proof}
For $a\in L^\times$, choose $F=K(a)$. Properness of the finite
extension gives $\int_{\Omega_F}\log|a|\,d\nu_F=0$.
The pullback integral formula gives the same value on $\Omega_L$.
Thus every nonzero element satisfies the product formula on $L$.
\end{proof}

\subsection{Compatibility in algebraic towers}

\begin{lemma}[Descent of a finite basis]\label{lem:finite-basis-descent}
\leanok
\source{chen-moriwaki-adelic:page-0213,chen-moriwaki-adelic:page-0214}
\uses{lem:finite-stage-exhaustion}
\lean{IntermediateField.exists_toSubmodule_eq_span_of_mul_mem,IntermediateField.exists_finite_basis_descent,IntermediateField.finiteDimensional_finrank_adjoin_of_basis_span}
Let $B/K$ be algebraic, let $N/B$ be finite, and let
$e_1,\ldots,e_n$ be a $B$-basis of $N$. Given finite subsets
$U\subseteq N$ and $T\subseteq B$, there is a finite subextension
$A/K$ in $B$ containing $T$ such that the $A$-span $E$ of the $e_i$
is a subfield of $N$ containing $U$. It satisfies
$[E:A]=[N:B]$ and $B(E)=N$.
More generally, the finite-dimensional span of a finite family over a
field $A$ is an intermediate field whenever it contains $1$ and the
products of all pairs from that family.
\end{lemma}
\begin{proof}
Adjoin to $K$ the set $T$, the coordinates of $1$, the coordinates of
every product $e_i e_j$, and the coordinates of every element of $U$.
These are finitely many elements of the algebraic extension $B/K$, so
they generate a finite extension $A/K$. Their defining coordinate
identities put $1$, all $e_i e_j$, and $U$ in the $A$-span $E$.
Bilinearity makes this span a subalgebra. In a finite-dimensional
subalgebra of a field, multiplication by a nonzero element is an
injective linear endomorphism and hence surjective. It therefore has
an inverse in the subalgebra, proving the general assertion as well.
The $e_i$ are $A$-linearly independent because they are
$B$-linearly independent; thus they are an $A$-basis of $E$ and
$[E:A]=n$. Their $B$-span is $N$, so $B(E)=N$.
\end{proof}

\begin{lemma}[A separable generator and purely inseparable powers]
\label{lem:separable-generator-powers}
\leanok
\source{chen-moriwaki-adelic:page-0213}
\lean{Field.exists_separable_generator_pow_aeval}
Let $N/B$ be finite and let $q$ be the exponential characteristic of
$B$. There is a separable element $c\in N$ whose minimal polynomial
has degree $[N:B]_s$, and every $x\in N$ satisfies
$x^{q^r}=P(c)$ for some $r\geq0$ and some $P\in B[T]$.
\end{lemma}
\begin{proof}
Let $C$ be the separable closure of $B$ in $N$. The primitive element
theorem gives $C=B(c)$ for a separable element $c$. Its minimal
polynomial has degree $[C:B]=[N:B]_s$. Since $N/C$ is purely
inseparable, every $x\in N$ has some $q$-power in $C$. Every element
of $C$ is a polynomial in its primitive element $c$, giving the
asserted expression. In characteristic zero, $q=1$ and $C=N$.
\end{proof}

\begin{lemma}[Finite descent with the separable part]\label{lem:finite-descent-separable}
\leanok
\source{chen-moriwaki-adelic:page-0213,chen-moriwaki-adelic:page-0214}
\uses{lem:finite-basis-descent,lem:separable-generator-powers}
\lean{IntermediateField.exists_finite_descent_of_finite_subset,IntermediateField.exists_finite_descent}
Let $B/K$ be algebraic, let $N/B$ be finite, and let $F/K$ be a
finite subextension of $N/K$. There are finite subextensions
$A/K$ in $B$ and $E/K$ in $N$, with $A,F\subseteq E$, such that
\[
 B(E)=N,\qquad [E:A]=[N:B],\qquad [E:A]_s=[N:B]_s.
\]
More generally, any prescribed finite subset of $N$ can be required
to lie in $E$.
\end{lemma}
\begin{proof}
Choose a $B$-basis $e_1,\ldots,e_n$ of $N$ and choose $c$ as in
Lemma~\ref{lem:separable-generator-powers}. Write
$e_i^{q^{r_i}}=P_i(c)$ with $P_i\in B[T]$.
Apply Lemma~\ref{lem:finite-basis-descent}, requiring the prescribed
finite subset and $c$ to belong to $E$, and requiring $A$ to contain
all coefficients of the minimal polynomial of $c$ over $B$ and of
the finitely many $P_i$. This gives the total-degree identity and
$B(E)=N$.

The minimal polynomial of $c$ over $A$ maps to its minimal polynomial
over $B$: the latter has coefficients in $A$ and remains irreducible
there, since a factorization over $A$ would give one over $B$.
Consequently $c$ is separable over $A$ and
$[A(c):A]=[N:B]_s$. Each identity $e_i^{q^{r_i}}=P_i(c)$ now has
coefficients in $A$, so every $e_i$ is purely inseparable over
$D=A(c)$. The elements of $E$ purely inseparable over $D$ form a
subfield containing $A$ and every $e_i$; hence they contain the
$A$-span $E$. Thus $E/D$ is purely inseparable, and the tower law
for separable degrees gives $[E:A]_s=[D:A]=[N:B]_s$.
To obtain $F\subseteq E$, prescribe a finite $K$-basis of $F$.
Since $K\subseteq A$, its $K$-span also lies in $E$.

These constructions also explain the scalar extensions: multiplication
maps $E\otimes_A B\to N$ isomorphically because it carries the
basis $e_i\otimes1$ to the $B$-basis $e_i$. Writing $C=B(c)$, the
same argument with the power basis of $c$ gives
$D\otimes_A B\simeq C$. Here $D$ and $C$ are the respective
separable closures, because the remaining extensions are purely
inseparable.
\end{proof}

\begin{lemma}[Restriction of embeddings after finite descent]
\label{lem:descended-embedding-bijection}
\leanok
\source{chen-moriwaki-adelic:page-0213,chen-moriwaki-adelic:page-0214}
\lean{IntermediateField.bijective_restrict_embeddings_of_adjoin_eq_top}
Let $A\subseteq B\subseteq N$ be fields, let $E$ be a subfield of $N$
containing $A$, and suppose that $E/A$ and $N/B$ are finite,
$B(E)=N$, and $[E:A]_s=[N:B]_s$. For every algebraically closed
extension $C$ of $B$, restriction is a bijection
\[
 \operatorname{Hom}_B(N,C)\ \longrightarrow\ \operatorname{Hom}_A(E,C).
\]
\end{lemma}
\begin{proof}
If two $B$-embeddings agree on $E$, then they agree on the field
generated by $B$ and $E$, which is $N$. Thus restriction is injective.
The two sets are finite, and their cardinalities are $[N:B]_s$ and
$[E:A]_s$, respectively. The assumed equality of these degrees makes
the injection a bijection.
\end{proof}

\begin{lemma}[Weighted place sums after finite descent]
\label{lem:descended-weighted-sums}
\leanok
\source{chen-moriwaki-adelic:page-0213,chen-moriwaki-adelic:page-0214}
\uses{lem:descended-embedding-bijection,lem:weighted-separable-embedding-sums,lem:algebraic-absolute-value-extension}
\lean{AbsoluteValue.sum_localFinSepDegree_nsmul_baseChange}
Under the hypotheses of Lemma~\ref{lem:descended-embedding-bijection},
fix an absolute value $y$ on $B$ and put $x=y|_A$. For an additive
commutative monoid $M$ and a function $\Phi$ from the absolute values
of $E$ to $M$,
\[
 \sum_{w|_B=y}[N_w:B_y]_s\,\Phi(w|_E)
 =
 \sum_{v|_A=x}[E_v:A_x]_s\,\Phi(v).
\]
\end{lemma}
\begin{proof}
Extend $y$ to an absolute value $u$ on an algebraic closure $C$ of $B$.
The weighted embedding formula changes the left sum into the sum of
$\Phi(u\circ\sigma|_E)$ over $\operatorname{Hom}_B(N,C)$ and the right
sum into the sum of $\Phi(u\circ\tau)$ over $\operatorname{Hom}_A(E,C)$.
The restriction bijection identifies these sums term by term.
\end{proof}

\begin{lemma}[Fibre probabilities under a descended base change]
\label{lem:descended-base-change}
\leanok
\source{chen-moriwaki-adelic:page-0213,chen-moriwaki-adelic:page-0214}
\uses{lem:descended-weighted-sums,lem:finite-absolute-fibre-sums}
\lean{ArakelovGeometry.AdelicStructure.ParameterSpaceFinExtension.baseChangeFiber,ArakelovGeometry.AdelicStructure.ParameterSpaceFinExtension.baseChangeFiber_ϕ,ArakelovGeometry.AdelicStructure.ParameterSpaceFinExtension.map_fibreProbabilityMeasure_baseChangeFiber,ArakelovGeometry.AdelicStructure.ParameterSpaceFinExtension.lintegral_baseChangeFiber}
Let $A\subseteq B\subseteq N$ and $E\subseteq N$ satisfy the hypotheses
of Lemma~\ref{lem:descended-embedding-bijection}. Let $S$ and $T$ be
adelic structures on $A$ and $B$, and let $\alpha\in\Omega_S$ and
$\beta\in\Omega_T$ satisfy $\phi_T(\beta)|_A=\phi_S(\alpha)$.
Restriction of absolute values defines
\[
 r:X_{N,\beta}\longrightarrow X_{E,\alpha}.
\]
It satisfies $|\cdot|_{r(w)}=|\cdot|_w|_E$ and
\[
 r_*P_{N,\beta}=P_{E,\alpha}.
\]
Consequently, for every $f:X_{E,\alpha}\to[0,\infty]$,
\[
 \int f(r(w))\,dP_{N,\beta}(w)=\int f(v)\,dP_{E,\alpha}(v).
\]
In particular this applies to the fields supplied by finite descent.
\end{lemma}
\begin{proof}
The matching restrictions to $A$ make $r$ well-defined. Each fibre is
a finite measurable space with all subsets measurable, so every map
between the two fibres is measurable. Its weights are the local
separable degrees divided by the global separable degrees.

Apply Lemma~\ref{lem:descended-weighted-sums} to an arbitrary
nonnegative function of the restricting absolute value, assigning it
value zero away from the extensions of $\phi_S(\alpha)$.
Divide by the common positive denominator $[E:A]_s=[N:B]_s$.
The finite weighted-sum formulas then identify its two integrals.
Taking indicator functions gives equality of the two probability
measures. The nonnegative change-of-variables formula gives the last
identity for every $f$; measurability is automatic on the finite fibre.
\end{proof}

\begin{remark}[The role of finite descent]\label{rem:algebraic-tower-source-repair}
\uses{lem:finite-descent-separable,lem:descended-base-change}
\source{chen-moriwaki-adelic:page-0213,chen-moriwaki-adelic:page-0214}
Finite descent replaces the prime-degree reduction in the printed
proof of Theorem 3.4.12. The assertion there that
$[F:F\cap B]$ divides $[N:B]$ is false without additional hypotheses:
in the splitting field of $T^3-2$ over $\mathbb Q$, let
$F=\mathbb Q(\sqrt[3]2)$ and
$B=\mathbb Q(\zeta_3\sqrt[3]2)$. These distinct cubic fields have
intersection $\mathbb Q$, whereas their compositum has degree two
over $B$. The descent lemma supplies the necessary equality of
degrees after enlarging the finite base, and also supplies the
equality of separable degrees needed for the probabilities.
\end{remark}

\begin{lemma}[Uniqueness on the projective parameter space]\label{lem:algebraic-total-measure-uniqueness}
\leanok
\uses{def:projective-measurable,def:finite-stages,lem:finite-stage-exhaustion,lem:projective-evaluations,lem:finite-stage-transitions-measurable}
\source{chen-moriwaki-adelic:page-0209,chen-moriwaki-adelic:page-0210}
\lean{ArakelovGeometry.AdelicStructure.Algebraic.totalMeasure_ext}
Let $\mu,\nu$ be measures on $(\Omega_L,\mathcal A_L)$ with $\mu$
finite. If $(\pi_{L/F})_*\mu=(\pi_{L/F})_*\nu$ for every finite
intermediate field $F/K$, then $\mu=\nu$.
\end{lemma}
\begin{proof}
Let $\mathcal C$ be the union of the pullback $\sigma$-algebras
$\pi_{L/F}^{-1}(\mathcal A_F)$. It contains the whole space and
generates $\mathcal A_L$. Two members represented at stages $F,G$
have their intersection represented at the finite compositum $FG$,
so $\mathcal C$ is a $\pi$-system. The pushforward hypothesis gives
equality of the two measures on $\mathcal C$, including equality of
their total masses. Both measures are therefore finite. Their
agreement sets form a Dynkin system by subtraction for complements
and countable additivity for disjoint unions. The $\pi$-$\lambda$
theorem gives agreement on all of $\mathcal A_L$.
\end{proof}

\begin{lemma}[Parameter and measurable-space compatibility]\label{lem:algebraic-tower-spaces}
\leanok
\lean{ArakelovGeometry.AdelicStructure.Algebraic.towerMeasurableEquiv,ArakelovGeometry.AdelicStructure.Algebraic.towerMeasurableEquiv_proj,ArakelovGeometry.AdelicStructure.Algebraic.towerMeasurableEquiv_ϕ,ArakelovGeometry.AdelicStructure.Algebraic.towerMeasurableEquiv_symm_proj}
\uses{def:algebraic-extension,prop:projective-measurable-equality,lem:restriction-transitive}
\source{chen-moriwaki-adelic:page-0212,chen-moriwaki-adelic:page-0214}
For $K\subseteq B\subseteq N$ algebraic, the iterated parameter
space and $\Omega_N$ identify by
$((\omega,v),w)\mapsto(\omega,w)$, with inverse
$(\omega,w)\mapsto((\omega,w|_B),w)$. This identification is a
measurable equivalence preserving absolute values and satisfying
$\pi_{N/K}=\pi_{B/K}\circ\pi_{N/B}$.
\end{lemma}
\begin{proof}
The restriction equalities show directly that the two displayed maps
are well defined and inverse. The direct $\sigma$-algebra is
generated by the base projection to $\Omega$ and evaluations on $N$.
The iterated one is generated by the projection to $\Omega_B$ and
evaluations on $N$. The projection to $\Omega_B$ is measurable for
the direct structure, since its base projection and its evaluations
at elements of $B$ are direct generators. Conversely, composing this
projection with $\pi_{B/K}$ gives the direct base projection.
Thus each generating family is measurable for the other structure.
The formulas for absolute values and projections follow from the maps.
\end{proof}

\begin{definition}[Fibre probabilities on the parameter space]\label{def:algebraic-total-fibre-probability}
\leanok
\uses{def:algebraic-extension,def:algebraic-fibre-probability,lem:algebraic-fibre-measurable,lem:algebraic-measurable-probability-family,def:algebraic-fibre-integral}
\source{chen-moriwaki-adelic:page-0210}
\lean{ArakelovGeometry.AdelicStructure.Algebraic.fibreMeasure,ArakelovGeometry.AdelicStructure.Algebraic.isProbabilityMeasure_fibreMeasure,ArakelovGeometry.AdelicStructure.Algebraic.measurable_fibreMeasure,ArakelovGeometry.AdelicStructure.Algebraic.lintegral_fibreMeasure}
Write $j_{L,\omega}:X_{L,\omega}\to\Omega_L$ for the inclusion and set
$\mu_{L,\omega}=(j_{L,\omega})_*P_{L,\omega}$. These are probability
measures, depend measurably on $\omega$, and for measurable
$f:\Omega_L\to[0,\infty]$ satisfy
$\int f\,d\mu_{L,\omega}=I^+_{L/K}f(\omega)$.
\end{definition}
\begin{proof}
The inclusion is measurable, so pushforward preserves the mass one.
Measurable dependence is the measurable-family assertion for fibre
probabilities. The pushforward integration formula gives
$\int f\,d\mu_{L,\omega}=\int(f\circ j_{L,\omega})\,dP_{L,\omega}$,
which is the definition of $I^+_{L/K}f(\omega)$.
\end{proof}

\begin{definition}[The iterated fibre probability]\label{def:algebraic-tower-probability}
\leanok
\uses{def:algebraic-total-fibre-probability,lem:algebraic-tower-spaces,prop:algebraic-kernel-measurable}
\source{chen-moriwaki-adelic:page-0213,chen-moriwaki-adelic:page-0214}
\lean{ArakelovGeometry.AdelicStructure.Algebraic.towerFibreMeasure,ArakelovGeometry.AdelicStructure.Algebraic.measurable_towerFibreKernel,ArakelovGeometry.AdelicStructure.Algebraic.isProbabilityMeasure_towerFibreMeasure,ArakelovGeometry.AdelicStructure.Algebraic.lintegral_towerFibreMeasure}
For $K\subseteq B\subseteq N$ algebraic, let $e$ identify the iterated
parameter space with $\Omega_N$. For each $\omega\in\Omega_K$ the formula
\[
 Q_\omega(A)=\int_{\Omega_B}
       (e_*\mu_{N,x})(A)\,d\mu_{B,\omega}(x)
       \qquad(A\in\mathcal A_N)
\]
defines a probability measure. The family $x\mapsto e_*\mu_{N,x}$ is
measurable, and every measurable $f:\Omega_N\to[0,\infty]$ satisfies
\[
 \int f\,dQ_\omega
 =I^+_{B/K}\bigl(I^+_{N/B}(f\circ e)\bigr)(\omega).
\]
\end{definition}
\begin{proof}
Pushforward along the measurable map $e$ preserves both measurable
dependence and total mass. Countable additivity of $Q_\omega$ follows
from monotone convergence applied to the sum of indicators of
pairwise disjoint sets. Its total mass is the integral of the
constant function one against a probability, hence is one.
The displayed integration formula first holds for indicators by the
definition, then for simple functions by finite additivity, and finally
for nonnegative measurable functions by increasing simple approximation.
The two pushforward integration formulas identify the resulting
iterated integral with the two fibre operators.
\end{proof}

\begin{lemma}[Finite upper stages containing a finite subextension]
\label{lem:algebraic-finite-upper-stage}
\leanok
\source{chen-moriwaki-adelic:page-0213,chen-moriwaki-adelic:page-0214}
\uses{lem:finite-stage-exhaustion}
\lean{ArakelovGeometry.AdelicStructure.Algebraic.exists_finite_upper_stage}
Let $L/K$ be algebraic, let $B$ be an intermediate field, and let
$H/K$ be a finite intermediate extension. There is an intermediate
field $N$ of $L/B$ that is finite over $B$ and contains $H$.
Equivalently, the inclusion $H\hookrightarrow L$ factors through a
$K$-embedding $H\hookrightarrow N$.
\end{lemma}
\begin{proof}
Choose a finite $K$-basis $e_1,\ldots,e_n$ of $H$ and put
$N=B(e_1,\ldots,e_n)$. Each $e_i$ is algebraic over $B$, since it is
algebraic over $K$. Adjoining finitely many algebraic elements gives
a finite extension of $B$. The field $N$ contains $K$ and each $e_i$,
so it contains their $K$-span $H$. The inclusion of $H$ into $N$
has the required composite with $N\hookrightarrow L$.
\end{proof}

\begin{lemma}[Algebraic tower integration on finite cylinders]
\label{lem:algebraic-tower-cylinders}
\leanok
\source{chen-moriwaki-adelic:page-0213,chen-moriwaki-adelic:page-0214}
\uses{lem:algebraic-finite-upper-stage,lem:finite-descent-separable,lem:descended-base-change,lem:algebraic-finite-embedding-projection,lem:finite-fibre-tower,lem:finite-fibre-embedding-naturality,lem:algebraic-tower-spaces}
\lean{ArakelovGeometry.AdelicStructure.Algebraic.fibreLIntegral_tower_projFinite}
Let $K\subseteq B\subseteq L$ be algebraic, let $H/K$ be finite in
$L$, and let $e:\Omega_{L/B/K}\to\Omega_{L/K}$ be the canonical
identification. For every function $g:\Omega_H\to[0,\infty]$ and
every $\omega\in\Omega_K$,
\[
 I^+_{B/K}\bigl(I^+_{L/B}(g\circ p_H\circ e)\bigr)(\omega)
 = I^+_{H/K}g(\omega).
\]
No global measurability assumption on $g$ is needed.
\end{lemma}
\begin{proof}
Choose a finite extension $N/B$ inside $L$ containing $H$.
Apply finite descent to the image of $H$ in $N$. It gives finite
$A/K$ in $B$ and finite $E/A$ in $N$, containing that image, with
$B(E)=N$ and $[E:A]_s=[N:B]_s$.
Write $j:H\hookrightarrow E$ for the resulting $K$-embedding,
$q:\Omega_{E/K}\to\Omega_{H/K}$ for restriction along $j$, and
$e_{E/A/K}$ for the finite tower identification. Put
$h=g\circ q\circ e_{E/A/K}$.

For a place $y\in\Omega_B$, let $x=p_A(y)$. Restriction maps the
finite fibre of $N/B$ above $y$ to the finite fibre of $E/A$ above
$x$, and preserves their probability measures by descended base
change. Restricting an $L$-place first to $N$, then to $E$, then
along $j$ gives its restriction to $H$. The algebraic finite
projection formula and this base-change formula therefore give
\[
 I^+_{L/B}(g\circ p_H\circ e)(y)=I^+_{E/A}h(p_A(y)).
\]
All test functions on the finite fibres are measurable, including
the restrictions of $g$, so the equality holds for arbitrary $g$.

The finite projection formula for $B/K$, followed by the finite
tower formula and naturality under the embedding $j$, now yields
\[
 \begin{aligned}
 I^+_{B/K}\bigl(I^+_{L/B}(g\circ p_H\circ e)\bigr)(\omega)
 &= I^+_{A/K}(I^+_{E/A}h)(\omega)\\
 &= I^+_{E/K}(g\circ q)(\omega)\\
 &= I^+_{H/K}g(\omega).
 \end{aligned}
\]
This proves the pointwise identity.
\end{proof}

\begin{lemma}[Signed integration against a total-space fibre probability]
\label{lem:algebraic-total-fibre-integrals}
\leanok
\source{chen-moriwaki-adelic:page-0211}
\uses{def:algebraic-total-fibre-probability,lem:algebraic-signed-integrability}
\lean{ArakelovGeometry.AdelicStructure.Algebraic.integral_fibreMeasure,ArakelovGeometry.AdelicStructure.Algebraic.integrable_fibreMeasure_ae}
Let $f:\Omega_L\to\mathbb R$ be measurable. Then
\[
 \int f\,d\mu_{L,\omega}=I_{L/K}f(\omega).
\]
If $f$ is $\nu_L$-integrable, it is $\mu_{L,\omega}$-integrable for
almost every $\omega$. As usual, a nonintegrable real function has
totalized integral zero in the displayed identity.
\end{lemma}
\begin{proof}
The measurable inclusion $j_{L,\omega}$ defining $\mu_{L,\omega}$
gives change of variables for $f^+$, $f^-$ and $|f|$. Therefore $f$
is integrable for the pushforward exactly when $f\circ j_{L,\omega}$
is integrable on the fibre. In the integrable case subtraction of the
two nonnegative formulas gives the stated signed identity; otherwise
both totalized integrals are zero. If $f$ is globally integrable,
nonnegative disintegration of $|f|$ shows that its integral on almost
every fibre is finite, giving the last assertion by this equivalence.
\end{proof}

\begin{lemma}[Signed integration against the iterated fibre probability]
\label{lem:algebraic-tower-probability-integral}
\leanok
\source{chen-moriwaki-adelic:page-0211,chen-moriwaki-adelic:page-0214}
\uses{def:algebraic-tower-probability,lem:algebraic-total-fibre-integrals}
\lean{ArakelovGeometry.AdelicStructure.Algebraic.integral_towerFibreMeasure}
In an algebraic tower $K\subseteq B\subseteq N$, let
$f:\Omega_N\to\mathbb R$ be measurable and $Q_\omega$-integrable.
Then
\[
 \int f\,dQ_\omega
 =I_{B/K}\bigl(I_{N/B}(f\circ e)\bigr)(\omega).
\]
\end{lemma}
\begin{proof}
Put $\eta_x=e_*\mu_{N,x}$, so
$Q_\omega(A)=\int\eta_x(A)\,d\mu_{B,\omega}(x)$.
The nonnegative integration formula applied to $|f|$ shows that
$f$ is $\eta_x$-integrable for $\mu_{B,\omega}$-almost every $x$,
and that these absolute integrals have finite integral. Applying the
same formula to $f^+$ and $f^-$ and subtracting gives
\[
 \int f\,dQ_\omega
 =\int_{\Omega_B}\left(\int f\,d\eta_x\right)d\mu_{B,\omega}(x).
\]
The inner integral is measurable: the nonnegative integrals of
$f^+$, $f^-$ and $|f|$ are measurable by the kernel formula; on the
measurable set where the last is finite it is the difference of the
first two, and elsewhere its totalized value is zero. Change of
variables along $e$ and Lemma~\ref{lem:algebraic-total-fibre-integrals}
identify it with $I_{N/B}(f\circ e)(x)$. A second application of
that lemma to the outer integral yields the asserted identity.
\end{proof}

\begin{lemma}[Equality of direct and iterated fibre probabilities]\label{lem:algebraic-tower-probability-equality}
\leanok
\uses{def:algebraic-tower-probability,lem:algebraic-tower-cylinders,lem:algebraic-total-measure-uniqueness,lem:algebraic-lintegral-pullback}
\source{chen-moriwaki-adelic:page-0213,chen-moriwaki-adelic:page-0214}
\lean{ArakelovGeometry.AdelicStructure.Algebraic.towerFibreMeasure_eq}
For every algebraic tower $K\subseteq B\subseteq N$ and
$\omega\in\Omega_K$, the iterated probability satisfies
$Q_\omega=\mu_{N,\omega}$ on $\mathcal A_N$.
\end{lemma}
\begin{proof}
For every finite $H/K$ in $N$ and every nonnegative measurable
$g$ on $\Omega_H$, the integral of $g\circ p_H$ against $Q_\omega$
is its iterated fibre integral. The cylinder tower identity
identifies this with $I^+_{H/K}g(\omega)$, which is also its
integral against $\mu_{N,\omega}$ by the finite marginal formula.
Taking indicators shows that the two probabilities have equal
pushforwards to every finite stage. Uniqueness on the projective
parameter space therefore gives the asserted equality.
\end{proof}

\begin{theorem}[Nonnegative integration in algebraic towers]\label{thm:algebraic-tower-lintegral}
\leanok
\uses{lem:algebraic-tower-probability-equality,def:algebraic-tower-probability,def:algebraic-total-fibre-probability,lem:algebraic-tower-spaces}
\source{chen-moriwaki-adelic:page-0212,chen-moriwaki-adelic:page-0214}
\lean{ArakelovGeometry.AdelicStructure.Algebraic.fibreLIntegral_tower}
For arbitrary algebraic $K\subseteq B\subseteq N$ and every
nonnegative measurable $g$ on the iterated parameter space,
\[
 I^+_{B/K}(I^+_{N/B}g)
 =I^+_{N/K}(g\circ e^{-1})
\]
pointwise on $\Omega_K$.
\end{theorem}
\begin{proof}
Put $f=g\circ e^{-1}$, which is measurable because $e$ is a
measurable equivalence. At each $\omega$, integrating $f$ against
$Q_\omega$ gives $I^+_{B/K}(I^+_{N/B}g)(\omega)$ by the iterated
integration formula and $e^{-1}\circ e=\mathrm{id}$. Equality
$Q_\omega=\mu_{N,\omega}$ identifies this with
$I^+_{N/K}(g\circ e^{-1})(\omega)$.
\end{proof}

\begin{corollary}[Tower integration with finite upper extension]\label{lem:algebraic-tower-finite-upper}
\leanok
\uses{thm:algebraic-tower-lintegral}
\source{chen-moriwaki-adelic:page-0213,chen-moriwaki-adelic:page-0214}
\lean{ArakelovGeometry.AdelicStructure.Algebraic.fibreLIntegral_tower}
If $B/K$ is algebraic and $N/B$ is finite, the same pointwise
nonnegative tower identity holds. In fact finiteness of $N/B$
may be replaced by algebraicity.
\end{corollary}
\begin{proof}
A finite extension is algebraic, so the preceding theorem applies
with its canonical identification of parameter spaces.
\end{proof}

\begin{lemma}[The equivalence of extended adelic structures]\label{lem:algebraic-tower-equivalence}
\leanok
\uses{thm:algebraic-tower-lintegral,lem:algebraic-tower-probability-equality,prop:algebraic-lintegral,lem:algebraic-tower-spaces,def:adelic-equivalence}
\source{chen-moriwaki-adelic:page-0212,chen-moriwaki-adelic:page-0214}
\lean{ArakelovGeometry.AdelicStructure.Algebraic.map_measure_towerMeasurableEquiv,ArakelovGeometry.AdelicStructure.Algebraic.towerEquiv}
The canonical parameter identification $e$ pushes the iterated
extension measure to the direct extension measure. Together with
its preservation of absolute values, it is an equivalence
$(S\otimes_KB)\otimes_BN\simeq S\otimes_KN$.
\end{lemma}
\begin{proof}
For a nonnegative measurable $f$ on $\Omega_N$, successive
disintegration writes its integral against the pushed-forward
iterated measure as
\[
 \int_{\Omega_K} I^+_{B/K}
       \bigl(I^+_{N/B}(f\circ e)\bigr)\,d\nu_K.
\]
The tower identity makes the integrand $I^+_{N/K}f$.
Disintegration for $N/K$ therefore gives $\int f\,d\nu_N$.
Indicators show equality of the measures. The parameter
identification already preserves absolute values, giving the
claimed equivalence.
\end{proof}

\begin{lemma}[Signed integration in an algebraic tower]\label{lem:algebraic-tower-signed}
\leanok
\uses{lem:algebraic-tower-equivalence,lem:algebraic-tower-probability-equality,lem:algebraic-total-fibre-integrals,lem:algebraic-tower-probability-integral,lem:algebraic-null-disintegration,lem:algebraic-fibre-integral-ae}
\source{chen-moriwaki-adelic:page-0211,chen-moriwaki-adelic:page-0214}
\lean{ArakelovGeometry.AdelicStructure.Algebraic.fibreIntegral_tower_of_measurable,ArakelovGeometry.AdelicStructure.Algebraic.fibreIntegral_tower}
For every integrable real function $g$ on the iterated parameter
space,
\[
 I_{B/K}(I_{N/B}g)=I_{N/K}(g\circ e^{-1})
 \quad\text{for }\nu_K\text{-almost every base point}.
\]
\end{lemma}
\begin{proof}
First suppose $g$ is measurable. The measure-preserving inverse of
$e$ makes $f=g\circ e^{-1}$ measurable and integrable on $\Omega_N$.
It is integrable for $\mu_{N,\omega}$ for almost every $\omega$.
At each such point, $Q_\omega=\mu_{N,\omega}$ makes $f$
$Q_\omega$-integrable. Signed integration against this iterated
probability identifies $\int f\,dQ_\omega$ with
$I_{B/K}(I_{N/B}g)(\omega)$, while integration against
$\mu_{N,\omega}$ identifies the same integral with
$I_{N/K}(g\circ e^{-1})(\omega)$.

For general integrable $g$, choose a measurable representative $g'$
equal to it almost everywhere. Almost everywhere congruence of fibre
integration, first for $N/B$ and then for $B/K$, makes the two
iterated integrals agree almost everywhere on $\Omega_K$.
Measure preservation transports $g=g'$ almost everywhere to the
direct space, and fibre congruence gives agreement of the direct
integrals almost everywhere. Applying the measurable case to $g'$
and combining these three equalities proves the assertion.
\end{proof}

\begin{theorem}[Compatibility of algebraic extensions]\label{thm:algebraic-tower}
\leanok
\lean{ArakelovGeometry.exists_equiv_algebraicExtension_algebraicExtension}
\uses{thm:algebraic-tower-lintegral,lem:algebraic-tower-spaces,lem:algebraic-tower-equivalence,lem:algebraic-tower-signed,prop:algebraic-lintegral,prop:algebraic-signed-disintegration,lem:algebraic-null-disintegration}
\source{chen-moriwaki-adelic:page-0212,chen-moriwaki-adelic:page-0214}
The measurable equivalence of Lemma~\ref{lem:algebraic-tower-spaces}
preserves the measures, and gives
$(S\otimes_K B)\otimes_B N\simeq S\otimes_K N$ as adelic structures.
It preserves the projections and satisfies the preceding pointwise
nonnegative integration formula. For any globally integrable real
function $f$ on either identified parameter space, one has
\[
 I_{B/K}(I_{N/B}f)=I_{N/K}f
 \quad\text{for }\nu\text{-almost every base point}.
\]
\end{theorem}
\begin{proof}
For a measurable set $A$, disintegrating the iterated measure gives
\[
 \nu_{N/B/K}(A)
 =\int_\Omega I^+_{B/K}(I^+_{N/B}\mathbf1_A)\,d\nu
 =\int_\Omega I^+_{N/K}\mathbf1_A\,d\nu=\nu_{N/K}(A).
\]
Thus the measurable equivalence is measure preserving, and its
absolute-value compatibility makes it an equivalence of adelic structures.

Suppose first $f$ is measurable and globally integrable. Apply the
nonnegative tower identity to $|f|$, $f^+$, and $f^-$. Their outer
integrals are finite, so outside a $\nu$-null set the direct
absolute-value fibre integral is finite and the iterated nonnegative
absolute-value integral is finite. At each remaining base point,
$f$ is integrable on the direct fibre, and is integrable on the
$N/B$ fibre for $P_{B,\omega}$-almost every intermediate point.
Moreover $I_{N/B}f$ is integrable under $P_{B,\omega}$, because its
absolute value is bounded almost everywhere by $I^+_{N/B}|f|$.
Subtracting the finite positive and negative tower identities proves
the signed identity at this base point.
For an arbitrary integrable representative, apply
Lemma~\ref{lem:algebraic-null-disintegration} successively to the
direct and iterated extensions. Changes on a global null set change
the inner integral only on a $\nu_B$-null set, and the outer integral
only on a $\nu$-null set. The measurable-representative identity
therefore gives the claimed almost everywhere identity.
\end{proof}

\begin{remark}[Domain of the signed tower formula]\label{rem:signed-tower-domain}
\uses{thm:algebraic-tower,prop:algebraic-integrability}
\source{chen-moriwaki-adelic:page-0210,chen-moriwaki-adelic:page-0211,chen-moriwaki-adelic:page-0213}
The tower formula is pointwise for nonnegative measurable functions,
and is an almost everywhere identity for globally integrable signed
functions. An unrestricted signed formula is not implied by assigning
zero to a nonintegrable real integral. The absolute-value criterion
therefore uses extended nonnegative integration, as does the proof of
the tower theorem in the source.
\end{remark}
